% =====================================================================
%  SRSI 2026/2027 — Projeto de Grupo, Parte 1 (teórica)
%  Template IEEE (IEEEtran, conferência, duas colunas) em português
%  Compilar: pdflatex -> bibtex -> pdflatex -> pdflatex  (ou latexmk -pdf)
%
%  Convenções do grupo
%  - \pendente{...}  = texto por escrever (sai a vermelho no PDF; apagar no fim)
%  - Cada secção tem um comentário "RESPONSÁVEL:" com a frente que a escreve
%  - Referências: só em referencias.bib (exportadas do Zotero, estilo IEEE)
%  - Figuras: pasta figuras/
% =====================================================================
\documentclass[conference,a4paper]{IEEEtran}

\usepackage[utf8]{inputenc}
\usepackage[T1]{fontenc}
\usepackage[portuguese]{babel}
\usepackage{cite}
\usepackage{amsmath,amssymb}
\usepackage{graphicx}
\usepackage{booktabs}
\usepackage{float}
\usepackage{multirow}
\usepackage{array}
\usepackage{xcolor}
\usepackage{tikz}
\usetikzlibrary{positioning,arrows.meta,fit}
\usepackage{pgfplots}
\pgfplotsset{compat=1.18}
\usepackage[hidelinks]{hyperref}
% Entrada do glossário (quebra entre colunas, ao contrário de uma tabela)
\newcommand{\gl}[2]{\par\noindent\hangindent=0.21\columnwidth\hangafter=1\makebox[0.21\columnwidth][l]{#1}\ignorespaces #2}

\renewcommand{\IEEEkeywordsname}{Palavras-chave}
\newcommand{\pendente}[1]{\textcolor{red}{[#1]}}
\newcommand{\checkok}{$\checkmark$} % célula "sim" nas tabelas comparativas
\newcommand{\nao}{--}            % célula "não"

\raggedbottom
% Tabelas grandes em páginas próprias, logo a seguir à página onde são referidas
\renewcommand{\dblfloatpagefraction}{0.25}
\renewcommand{\floatpagefraction}{0.25}
\begin{document}

% ---------------------------------------------------------------------
%  CAPA (dados pedidos no template do professor)
% ---------------------------------------------------------------------
\title{
  \includegraphics[width=0.42\textwidth]{figuras/logo-iscte.png}\\[6pt]
  {\normalsize Mestrado em Engenharia de Telecomunicações e Informática\\
   Segurança em Redes e Sistemas de Informação\\[6pt]}
  Deteção e Contenção Automática de Dispositivos IoT Comprometidos com
  \textit{Machine Learning}: Uma Revisão Sistemática da Literatura
}

\author{
\IEEEauthorblockN{Afonso Duarte Rosado Lopes (122631),
Rafael Fernandes dos Reis (122616),
Rafael Andrade da Silva (122638),\\
Gonçalo José Bilé Batista (122623),
João Azinheira Alves (123488)}
\IEEEauthorblockA{Iscte -- Instituto Universitário de Lisboa\\
Docentes: Prof. Valderi Reis Quietinho Leithardt e Prof. Maurício Machado da Rosa\\
Data: 29/10/2026}
}

\maketitle

% ---------------------------------------------------------------------
%  RESUMO  (máx. meia página, sem referências nem links)
%  RESPONSÁVEL: Frente 1 (escrever no fim)
% ---------------------------------------------------------------------
\begin{abstract}
\addcontentsline{toc}{section}{Resumo}
Os dispositivos da Internet das Coisas (IoT) têm poucos recursos, raramente são
atualizados e são facilmente comprometidos por \textit{malware} que os recruta para
\textit{botnets}, como a Mirai. Quando um deles é comprometido,
continua a funcionar normalmente aos olhos do utilizador, e o tráfego de rede é
muitas vezes a única forma de o detetar. Este trabalho apresenta uma revisão
sistemática da literatura sobre a deteção de dispositivos IoT comprometidos com
\textit{machine learning} a partir do tráfego de rede, com foco no passo seguinte à
deteção: a resposta automática que contém o dispositivo. Seguindo o protocolo
da unidade curricular, foram pesquisadas cinco bases de dados (677
resultados). Os estudos foram selecionados com critérios de inclusão e exclusão, um
critério específico sobre a existência de resposta automática e uma avaliação de
qualidade, resultando em 31 estudos incluídos. A análise mostra que muitos
trabalhos anunciam mecanismos de mitigação mas poucos os implementam e avaliam, e que
a avaliação em protótipo ou em tempo real continua a ser minoritária. Com base nestas
conclusões, propõe-se uma arquitetura para um \textit{gateway} que extrai fluxos com
Zeek, compara modelos supervisionados e não supervisionados e coloca automaticamente
em quarentena os dispositivos suspeitos, enquadrada numa arquitetura \textit{Zero
Trust} segundo o NIST SP 800-207.
\end{abstract}

\begin{IEEEkeywords}
Internet das Coisas, deteção de intrusões, \textit{machine learning}, \textit{botnets}, contenção automática, revisão sistemática
\end{IEEEkeywords}

% ---------------------------------------------------------------------
%  GLOSSÁRIO (obrigatório no template do professor)
%  RESPONSÁVEL: todos acrescentam as siglas que usam
% ---------------------------------------------------------------------
% ÍNDICE (pedido no template do professor)
\renewcommand{\contentsname}{Índice}
\setcounter{tocdepth}{1}
\noindent\begin{minipage}{\columnwidth}\footnotesize
\makeatletter\renewcommand{\l@section}[2]{\@dottedtocline{1}{0pt}{2.2em}{#1}{#2}}\makeatother
\tableofcontents
\end{minipage}

\newpage
\section*{Glossário}
\addcontentsline{toc}{section}{Glossário}
{\footnotesize
\gl{ACL}{Lista de controlo de acessos (\textit{Access Control List})}
\gl{AE}{\textit{Autoencoder}}
\gl{AME}{Eficácia de mitigação de ataques (\textit{Attack Mitigation Effectiveness})}
\gl{AUC, AUROC}{Área sob a curva ROC (\textit{Area Under the ROC Curve})}
\gl{AVAC}{Aquecimento, ventilação e ar condicionado}
\gl{C\&C}{Comando e controlo (\textit{Command and Control})}
\gl{CIL, DIL}{\textit{Class Incremental Learning}, \textit{Domain Incremental Learning}}
\gl{CISA}{\textit{Cybersecurity and Infrastructure Security Agency}}
\gl{CNN}{Rede neuronal convolucional (\textit{Convolutional Neural Network})}
\gl{CoAP}{\textit{Constrained Application Protocol}}
\gl{CPS}{Sistema ciber-físico (\textit{Cyber-Physical System})}
\gl{DDoS}{Negação de serviço distribuída (\textit{Distributed Denial of Service})}
\gl{DL}{Aprendizagem profunda (\textit{Deep Learning})}
\gl{DNS}{\textit{Domain Name System}}
\gl{DoS}{Negação de serviço (\textit{Denial of Service})}
\gl{DT}{Árvore de decisão (\textit{Decision Tree})}
\gl{EC}{Critério de exclusão (\textit{Exclusion Criterion})}
\gl{FAR}{Taxa de falsos alarmes (\textit{False Alarm Rate})}
\gl{FN, FP}{Falsos negativos, falsos positivos}
\gl{FNR}{Taxa de falsos negativos (\textit{False Negative Rate})}
\gl{FPR}{Taxa de falsos positivos (\textit{False Positive Rate})}
\gl{GNN}{Rede neuronal de grafos (\textit{Graph Neural Network})}
\gl{GRU, BiGRU}{\textit{Gated Recurrent Unit} (bidirecional)}
\gl{IC}{Critério de inclusão (\textit{Inclusion Criterion})}
\gl{IDS}{Sistema de deteção de intrusões (\textit{Intrusion Detection System})}
\gl{IF}{\textit{Isolation Forest}}
\gl{IoT}{Internet das Coisas (\textit{Internet of Things})}
\gl{IP}{\textit{Internet Protocol}}
\gl{IPS, IRS}{Sistema de prevenção (\textit{Intrusion Prevention System}) e de resposta a intrusões (\textit{Intrusion Response System})}
\gl{IPFIX}{\textit{IP Flow Information Export}}
\gl{KNN}{\textit{k-Nearest Neighbors}}
\gl{LR}{Regressão logística (\textit{Logistic Regression})}
\gl{LSTM}{\textit{Long Short-Term Memory}}
\gl{MAC}{\textit{Media Access Control}}
\gl{MCC}{\textit{Matthews Correlation Coefficient}}
\gl{MFA}{Autenticação multifator (\textit{Multi-Factor Authentication})}
\gl{MITM}{\textit{Man-in-the-Middle}}
\gl{ML}{Aprendizagem automática (\textit{Machine Learning})}
\gl{MLP}{\textit{Multilayer Perceptron}}
\gl{MQTT}{\textit{Message Queuing Telemetry Transport}}
\gl{MUD}{\textit{Manufacturer Usage Description}}
\gl{NB}{Naive Bayes}
\gl{NIST}{\textit{National Institute of Standards and Technology}}
\gl{P4}{\textit{Programming Protocol-independent Packet Processors}}
\gl{PA}{\textit{Policy Administrator}}
\gl{PAM}{\textit{Polygon Area Metric}}
\gl{PDP}{\textit{Policy Decision Point}}
\gl{PDR}{Taxa de entrega de pacotes (\textit{Packet Delivery Ratio})}
\gl{PE}{\textit{Policy Engine}}
\gl{PEP}{\textit{Policy Enforcement Point}}
\gl{PIP}{\textit{Policy Information Point}}
\gl{PIPOC}{População, Intervenção, Comparação, Resultados (\textit{Outcomes}), Contexto}
\gl{QA}{Avaliação de qualidade (\textit{Quality Assessment})}
\gl{QE, QF, QG}{Questão estatística, questão focada, questão geral}
\gl{RF}{\textit{Random Forest}}
\gl{RGPD}{Regulamento Geral sobre a Proteção de Dados}
\gl{RL}{Aprendizagem por reforço (\textit{Reinforcement Learning})}
\gl{RNN}{Rede neuronal recorrente (\textit{Recurrent Neural Network})}
\gl{RPL}{\textit{Routing Protocol for Low-Power and Lossy Networks}}
\gl{RSL}{Revisão sistemática da literatura}
\gl{SDN}{Redes definidas por \textit{software} (\textit{Software-Defined Networking})}
\gl{SP}{\textit{Special Publication} (série de documentos do NIST)}
\gl{SQLi}{Injeção de SQL (\textit{SQL Injection})}
\gl{SRSI}{Segurança em Redes e Sistemas de Informação}
\gl{SVM}{\textit{Support Vector Machine}}
\gl{TCP, UDP}{\textit{Transmission Control Protocol}; \textit{User Datagram Protocol}}
\gl{TN, TP}{Verdadeiros negativos, verdadeiros positivos}
\gl{TPR}{Taxa de verdadeiros positivos (\textit{True Positive Rate}), igual ao \textit{recall}}
\gl{VLAN}{\textit{Virtual Local Area Network}}
\gl{WSN}{Rede de sensores sem fios (\textit{Wireless Sensor Network})}
\gl{XAI}{Inteligência artificial explicável (\textit{eXplainable Artificial Intelligence})}
\gl{XGBoost}{\textit{Extreme Gradient Boosting}}
\gl{XSS}{\textit{Cross-Site Scripting}}
\gl{ZTA}{Arquitetura \textit{Zero Trust} (\textit{Zero Trust Architecture})}
\gl{ZTMM}{\textit{Zero Trust Maturity Model} (CISA)}
\par}

% =====================================================================
\section{Introdução}
% RESPONSÁVEL: Frente 1 (Afonso)
% =====================================================================
A Internet das Coisas (IoT) deixou de ser um conceito de nicho e passou a fazer parte
de casas, hospitais, fábricas e cidades. Câmaras de vigilância, termóstatos, fechaduras,
sensores e controladores industriais estão hoje ligados à rede de forma permanente e,
em muitos casos, diretamente expostos à Internet. A maioria destes equipamentos foi
pensada para ser barata e fácil de instalar, e não para resistir a ataques. Têm pouca
memória e pouca capacidade de processamento, raramente recebem atualizações de
\textit{firmware} e não têm margem para correr um antivírus ou um agente de segurança.
Muitos continuam a usar as credenciais de fábrica durante toda a vida útil.

A \textit{botnet} Mirai mostrou em 2016 o que acontece quando estas fragilidades são
exploradas em grande escala. O \textit{malware} percorria a Internet à procura de
dispositivos com o serviço Telnet acessível e tentava entrar com uma pequena lista de
credenciais por omissão. No pico da sua atividade controlava cerca de 600 mil
equipamentos, que foram usados em ataques de negação de serviço distribuída (DDoS)
contra alvos como o fornecedor de DNS Dyn~\cite{antonakakis2017mirai}. O código-fonte
acabou por ser publicado e deu origem a inúmeras variantes, mas o padrão mantém-se: aos
olhos do dono, o dispositivo infetado continua a funcionar normalmente, enquanto em
segundo plano faz varrimentos, participa em ataques ou comunica com um servidor de
comando e controlo (C\&C).

Como não é possível instalar \textit{software} de proteção nestes dispositivos, quem
gere a rede depende quase só do tráfego para perceber que algo mudou. Esse tráfego pode
ser observado num ponto central, como o \textit{gateway}, sem tocar nos equipamentos. É
neste cenário que a aprendizagem automática (ML) tem sido mais explorada. A partir de
estatísticas de fluxos, como a duração, o volume, os intervalos entre pacotes ou as
portas contactadas, os modelos aprendem a separar o comportamento habitual de cada
equipamento dos padrões típicos de uma infeção. A literatura sobre esta forma de deteção
é extensa e reporta com frequência taxas de acerto acima de 99\% em \textit{datasets}
públicos, como se vê em vários dos estudos incluídos nesta
revisão~\cite{ravi2020semisupervised,elsayed2024mpguard,toony2024multiblock,gowtham2026fedmp,Iqbal2026}.

Detetar o problema, porém, é só metade do trabalho. Um alerta que fica à espera de ser
lido deixa o dispositivo ativo na rede durante minutos ou horas, tempo suficiente para
participar num ataque ou para infetar outros equipamentos. Numa casa ou numa pequena
empresa, muitas vezes nem existe alguém para o ler. Faz por isso sentido que a deteção
seja acompanhada de uma resposta automática, aplicada no mesmo ponto onde o tráfego é
observado: bloquear o dispositivo, limitar o seu débito ou isolá-lo numa rede de
quarentena. O trabalho de Noetzold \textit{et al.}~\cite{noetzold2026anomaly}, que serve
de ponto de partida a este projeto, integra modelos de \textit{deep learning} num sistema
de monitorização que deteta anomalias em vários domínios, entre eles processos suspeitos
a partir de fluxos do CICIDS2017, e apresenta os alertas numa interface em tempo real.
Entre as direções de trabalho futuro, os autores apontam a integração com sistemas de
resposta a incidentes, a execução na \textit{edge} e o uso de aprendizagem não
supervisionada, e é precisamente esse passo que aqui se estuda.

Existem já várias revisões sobre segurança em IoT. Negera \textit{et al.}~\cite{negera2022review}
analisam a deteção de \textit{botnets} com ML em redes IoT geridas por SDN; Kandah
\textit{et al.}~\cite{kandah2025navigating} percorrem arquiteturas, ameaças emergentes e
contramedidas adaptativas; Vera-Mundaca e Aquino-Trujillo~\cite{veramundaca2024cybersecurity}
concentram-se nas ameaças e nas formas de mitigação em casas inteligentes. São visões
abrangentes, mas nenhuma tem como objetivo central verificar se os mecanismos de
contenção propostos chegam a ser implementados e avaliados, nem onde e como atuam. Esta
revisão foca-se nesse ponto.

\textbf{Objetivos.} O trabalho tem dois objetivos. O primeiro é perceber, através de uma
revisão sistemática da literatura (RSL), como é que a análise de tráfego com ML tem sido
usada para detetar e conter automaticamente dispositivos IoT comprometidos. Interessa
saber que modelos e \textit{features} são usados, com que \textit{datasets} e métricas são
avaliados, que mecanismos de resposta são propostos e onde atuam, e se as soluções foram
testadas para lá de um \textit{dataset} estático. Estas questões estão formalizadas nas
perguntas de pesquisa da Tabela~\ref{tab:perguntas}. O segundo objetivo é usar o
resultado da revisão para desenhar a solução que será implementada na Parte~2 do
projeto: um \textit{gateway} que analisa o tráfego dos dispositivos e, ao detetar um
comprometimento, os coloca automaticamente em quarentena.

\textbf{Contribuições.} As contribuições deste trabalho são quatro:
\begin{itemize}
  \item uma RSL de 31 estudos publicados entre 2020 e 2026, selecionados a
  partir de 677 resultados em cinco bases de dados, com um critério explícito sobre a
  existência de resposta automática e uma avaliação de qualidade de cada estudo, que
  mostra que só 14 deles descrevem e avaliam a resposta e só sete medem o tempo que esta
  demora a conter um ataque;
  \item uma comparação dos trabalhos relacionados (Tabelas~\ref{tab:comparacao}
  e~\ref{tab:comparacao2}) centrada
  na forma como a resposta é desencadeada, onde é aplicada e como é avaliada;
  \item a proposta de uma arquitetura de deteção e contenção para a Parte~2, assente na
  extração de fluxos com Zeek, em modelos supervisionados e não supervisionados e na
  quarentena com nftables, a avaliar com o \textit{dataset} IoT-23~\cite{iot23};
  \item o enquadramento dessa arquitetura no modelo \textit{Zero Trust} do NIST
  SP 800-207 e no modelo de maturidade da CISA, aplicado a um edifício inteligente.
\end{itemize}

\textbf{Estrutura.} A Secção~\ref{sec:fundamentos} apresenta os conceitos base.
A Secção~\ref{sec:metodologia} descreve o protocolo da revisão sistemática.
A Secção~\ref{sec:resultados} responde às perguntas de pesquisa e compara os
trabalhos relacionados. A Secção~\ref{sec:proposta} apresenta a solução proposta e a
Secção~\ref{sec:zta} enquadra-a numa arquitetura \textit{Zero Trust}. A
Secção~\ref{sec:discussao} discute as lacunas e as ameaças à validade e a
Secção~\ref{sec:conclusoes} conclui.

% =====================================================================
\section{Fundamentos}\label{sec:fundamentos}
% Secção curta: só o necessário para perceber o resto.
% =====================================================================
\subsection{Ameaças em dispositivos IoT}
% RESPONSÁVEL: Frente 2 (João Alves)
Um dispositivo IoT comprometido é ao mesmo tempo vítima e plataforma de ataque, e por isso
não basta detetá-lo: é preciso contê-lo. Três características tornam estes dispositivos
vulneráveis. A primeira são os recursos limitados: com pouca capacidade de processamento, pouca
memória e pouca energia, não conseguem correr um antivírus nem cifra pesada, e a defesa tende a subir para o
\textit{gateway}, para o \textit{edge} ou para o controlador SDN~\cite{kponyo2020lightweight}.
A segunda é a segurança fraca por omissão: credenciais por defeito, \textit{firmware}
raramente atualizado e protocolos leves, como o MQTT, o CoAP, o Zigbee ou o LoRaWAN, em que
a autenticação e a cifra não são obrigatórias~\cite{mqtt5,rfc7252}. Foi das credenciais
por omissão que a Mirai se aproveitou para se propagar~\cite{antonakakis2017mirai}. A
terceira é a heterogeneidade e a escala: milhares de dispositivos de fabricantes e
protocolos diferentes tornam difícil manter uma política de segurança uniforme.

A Tabela~\ref{tab:ameacas} resume as famílias de ataque consideradas nesta revisão e o
sinal que cada uma deixa no tráfego. Um dispositivo comprometido altera o seu
comportamento de rede: abre mais fluxos, contacta destinos novos, muda de ritmo. É este
desvio, observável no tráfego, que os modelos de ML aprendem a identificar.

\begin{table}[H]
\caption{Famílias de ataque a dispositivos IoT consideradas na revisão}\label{tab:ameacas}
\centering\scriptsize
\setlength{\tabcolsep}{3pt}
\begin{tabular}{@{}>{\raggedright\arraybackslash}p{0.21\columnwidth}>{\raggedright\arraybackslash}p{0.38\columnwidth}>{\raggedright\arraybackslash}p{0.35\columnwidth}@{}}
\toprule
Família & O que é & Sinal no tráfego \\
\midrule
DoS/DDoS & Inundação de um alvo (SYN, UDP, ICMP, HTTP) até esgotar a largura de banda ou os recursos & Pico de pacotes ou fluxos por segundo; muitas origens para um destino \\
\textit{Botnet} e C\&C & Dispositivos infetados que recebem ordens de um servidor de comando e controlo e propagam a infeção, como na Mirai~\cite{antonakakis2017mirai} & Ligações periódicas a um destino externo; varrimento de portas (Telnet); tráfego de saída anómalo \\
Força bruta & Tentativas repetidas de \textit{login} ou de autenticação & Muitas ligações curtas ao mesmo serviço \\
\textit{Malware} e injeção & Código malicioso, injeção de comandos ou de dados, exploração de vulnerabilidades & \textit{Payloads} e padrões de protocolo fora do normal \\
Reconhecimento & Varrimento de portas e \textit{fingerprinting} que preparam o ataque & Muitos destinos ou portas, com poucos \textit{bytes} por fluxo \\
Ataques de protocolo e de rede & MITM, \textit{spoofing}, \textit{black-hole}, \textit{selective forwarding} (WSN), ataques ao MQTT~\cite{karam2025isaaf} & Anomalias de endereçamento, encaminhamento ou temporização \\
\bottomrule
\end{tabular}
\end{table}

\subsection{Deteção de intrusões baseada em fluxos}
% RESPONSÁVEL: Frente 2 (João Alves)
A deteção por fluxos analisa o resumo estatístico de uma comunicação, e não o conteúdo dos
pacotes, o que a torna leve e compatível com tráfego cifrado. Um fluxo é o conjunto de
pacotes que partilham o mesmo 5-tuplo (endereços IP e portas de origem e de destino, e
protocolo) dentro de uma janela de tempo ou até expirar um
\textit{timeout}~\cite{rfc7011}. Os fluxos podem ser exportados pelos próprios
equipamentos de rede, com NetFlow ou IPFIX~\cite{rfc3954,rfc7011}, extraídos de capturas
com ferramentas como o CICFlowMeter~\cite{lashkari2017cicflowmeter} ou o
Zeek~\cite{paxson1999bro}, ou lidos pelo controlador SDN a partir dos contadores dos
\textit{switches} OpenFlow.

O processo típico tem quatro passos: captar o tráfego no \textit{gateway}, no
\textit{edge} ou no \textit{switch}; agregá-lo em fluxos ou em janelas de tempo; calcular
as \textit{features}; e classificá-lo com um modelo de ML, que decide se o fluxo ou o
dispositivo é benigno ou malicioso. A Tabela~\ref{tab:features} apresenta os grupos de
\textit{features} mais comuns.

\begin{table}[H]
\caption{Grupos de features habitualmente extraídas de fluxos de rede}\label{tab:features}
\centering\scriptsize
\begin{tabular}{@{}>{\raggedright\arraybackslash}p{0.30\columnwidth}>{\raggedright\arraybackslash}p{0.64\columnwidth}@{}}
\toprule
Grupo & Exemplos \\
\midrule
Volume & Número de pacotes e de \textit{bytes} por fluxo; pacotes e \textit{bytes} por segundo \\
Temporais & Duração do fluxo; tempo entre pacotes (média, desvio-padrão) \\
Tamanho dos pacotes & Mínimo, máximo, média e desvio-padrão \\
\textit{Flags} e protocolo & Contagem de \textit{flags} TCP (SYN, ACK, RST); protocolo; portas \\
Agregadas por dispositivo ou janela & Número de destinos distintos e de fluxos por origem; razão entre tráfego de entrada e de saída \\
\bottomrule
\end{tabular}
\end{table}

A abordagem tem vantagens claras. Como não inspeciona o \textit{payload}, custa menos e
expõe menos dados pessoais, funciona com tráfego cifrado, escala bem e serve para
dispositivos que não podem correr \textit{software} de segurança. Em SDN, os contadores de
fluxo já existem no controlador, o que facilita o seu uso. Tem também limites. Sem
\textit{payload}, a injeção de comandos e o \textit{malware} escondido em tráfego
aparentemente normal são difíceis de detetar. A agregação em janelas introduz atraso: no
R-IDPS, por exemplo, a deteção opera sobre janelas de 10\,s~\cite{mazhar2022ridps}. E os
modelos dependem de \textit{datasets} rotulados, muitas vezes simulados, o que afeta a
generalização, como se mostra na QF2 (Secção~\ref{sec:qf2}).

\subsection{Aprendizagem automática para deteção}
% RESPONSÁVEL: Frentes 3 (Rafael Silva, supervisionado) e 4 (Rafael Reis, não supervisionado)
Com o aumento de dispositivos IoT em casas e empresas, a exposição a todo o tipo de
ataques tem vindo a crescer. Estes ataques podem ser conhecidos, com diferenças já
identificadas entre o seu tráfego e o tráfego normal da rede, ou podem ser totalmente
novos, sem nenhum rótulo de ataque conhecido.

Devido a esta diferença nos tipos de ataque que podem acontecer, os modelos de ML também
diferem, sobretudo na forma como detetam estes ataques e na sua eficácia para cada um
deles.

Os modelos que utilizam os dados já conhecidos sobre os vários ataques têm uma
aprendizagem supervisionada. Na fase de treino, o modelo já sabe que tráfego é malicioso
e procura os padrões que esse tráfego tem em comum, para depois identificar e rotular
os ataques em tráfego novo. Um exemplo deste tipo de aprendizagem é a \textit{Random
Forest}.

\subsubsection{\textit{Random Forest}}
A \textit{Random Forest} (RF)~\cite{breiman2001random} é um algoritmo de classificação
bastante usado na deteção de ataques, que permite também identificar as
\textit{features} de tráfego mais importantes para os detetar. Com isso, é possível
reduzir o custo computacional da deteção, já que deixa de ser necessário analisar todas
as \textit{features}, mas apenas as indicadas~\cite{Iranzad2025}.

O algoritmo envolve a construção de uma coleção de árvores de decisão, cada uma treinada
com uma amostra aleatória dos dados, que classificam o tráfego como benigno ou
malicioso; a classe final é a escolhida pela maioria das árvores. As árvores integram
também a ideia de \textit{random node split}: em cada nó, apenas um subconjunto
aleatório das \textit{features} de tráfego é considerado.

Com esta ideia aplicada, a RF consegue medir a importância de cada \textit{feature}. Para
cada árvore, calcula-se a exatidão (\textit{accuracy}) nos dados que ficaram fora da sua
amostra e volta a calcular-se depois de embaralhar, aleatoriamente, os valores de uma das
\textit{features}. Se a exatidão cair muito com os novos valores, a \textit{feature}
alterada é importante para distinguir o tráfego malicioso e deve ser mantida, para que o
modelo detete os ataques com mais sucesso~\cite{breiman2001random,Iranzad2025}.

Apesar de ser muito bom a identificar ataques já conhecidos, este modelo é
consideravelmente pior na deteção de ataques novos, para os quais não existem exemplos
nem rótulos no treino. Como muitas vezes não existem dados rotulados em tempo real e os
ataques mudam constantemente, os modelos que não dependem de rótulos de ataque tornam-se
importantes: aprendem o comportamento normal da rede e assinalam o que se afasta dele.
Dois são usados neste trabalho.

\subsubsection{\textit{Isolation Forest}}
O \textit{Isolation Forest} (iForest)~\cite{liu2008isolation} é um algoritmo não
supervisionado eficiente para identificar anomalias em grandes volumes de tráfego. Ao
contrário dos métodos baseados em distâncias ou densidades, aproveita duas
características próprias das anomalias: são raras e têm valores muito diferentes do
comportamento habitual, o que facilita o seu isolamento.

O algoritmo constrói um conjunto de árvores binárias, as \textit{isolation trees}. Em
cada nó escolhe aleatoriamente uma \textit{feature} e um valor de corte entre o mínimo e
o máximo dessa \textit{feature}. Como as anomalias precisam de menos partições para
ficarem isoladas, aparecem mais perto da raiz, com caminhos mais curtos. Seja
$X=\{x_1,\dots,x_n\}$ um conjunto com $n$ amostras e $h(x)$ o comprimento do caminho de
$x$, isto é, o número de ligações da raiz até à folha onde $x$ termina. O \textit{score}
de anomalia é
\begin{equation}
s(x,n) = 2^{-\frac{E[h(x)]}{c(n)}},\qquad c(n) = 2H(n-1) - \frac{2(n-1)}{n}
\label{eq:iforest}
\end{equation}
em que $E[h(x)]$ é a média de $h(x)$ em todas as árvores e $c(n)$ é o comprimento médio
de caminho numa árvore binária com $n$ amostras, que serve para normalizar. O número
harmónico é aproximado por $H(i)\approx \ln(i)+\gamma$, com $\gamma \approx 0{,}5772$ (a
constante de Euler-Mascheroni). O \textit{score} varia em $(0,1]$: valores próximos de 1
indicam uma amostra isolada com muita facilidade, provavelmente um ataque; valores
claramente abaixo de 0,5 indicam tráfego normal; e se todas as amostras tiverem
$s\approx 0{,}5$, o conjunto não tem anomalias distintas~\cite{liu2008isolation}.

\subsubsection{\textit{Autoencoder}}
O \textit{autoencoder} é uma rede neuronal simétrica treinada para reproduzir a própria
entrada. Quando é treinado só com tráfego benigno, aprende a reconstruí-lo bem, e falha
na reconstrução de padrões que nunca viu, como os de um DDoS ou de uma
\textit{botnet}~\cite{mirsky2018kitsune}. Na literatura é muitas vezes classificado como
não supervisionado, porque não usa rótulos de ataque; neste trabalho é tratado como
semi-supervisionado, porque o treino exige saber que o tráfego usado é benigno.

A rede tem duas partes. O codificador comprime a entrada $x\in\mathbb{R}^d$, com as
\textit{features} de um fluxo, numa camada oculta $h\in\mathbb{R}^p$ com $p<d$:
\begin{equation}
h = f(W_1 x + b_1)
\label{eq:ae_enc}
\end{equation}
em que $W_1$ é a matriz de pesos, $b_1$ o vetor de \textit{bias} e $f(\cdot)$ uma função
de ativação não linear, como a ReLU ou a sigmoide. O descodificador tenta reconstruir a
entrada a partir dessa representação:
\begin{equation}
\hat{x} = g(W_2 h + b_2)
\label{eq:ae_dec}
\end{equation}
O desvio entre a entrada e a reconstrução mede-se, para cada amostra, pelo erro
quadrático médio. A rede é treinada para minimizar a média deste erro nas $N$ amostras
de treino, e uma amostra é classificada como anómala quando o seu erro ultrapassa um
limiar $\theta$, definido a partir do erro observado no tráfego benigno de validação:
\begin{equation}
e(x) = \frac{1}{d}\sum_{j=1}^{d}\left(x_j-\hat{x}_j\right)^2,\qquad
\text{anomalia} \iff e(x) > \theta
\label{eq:ae}
\end{equation}

A diferença entre os dois modelos tem consequências práticas. O \textit{Isolation
Forest} tem custo computacional baixo e não precisa de dados de treino rotulados, o que
o torna adequado para equipamentos com poucos recursos. O \textit{autoencoder} exige mais
capacidade de treino, mas consegue captar relações não lineares entre as
\textit{features}, e é por isso mais usado em análise centralizada de
fluxos~\cite{mirsky2018kitsune,zolotukhin2021curious}.

\subsection{Mecanismos de resposta e contenção}
% RESPONSÁVEL: Frente 5 (Gonçalo Batista)
A deteção de uma intrusão não garante, por si só, que a ameaça seja interrompida. Depois de identificado um comportamento malicioso, é necessário aplicar medidas que impeçam o dispositivo comprometido de continuar a atacar outros equipamentos ou a comprometer o funcionamento da rede. Enquanto um sistema de deteção de intrusões (IDS) pode limitar-se a gerar alertas, um sistema de prevenção (IPS) ou de resposta a intrusões (IRS) procura atuar diretamente sobre a ameaça. Em ambientes IoT, onde nem sempre existe um administrador disponível para intervir, a automatização deste processo torna-se particularmente importante~\cite{kponyo2020lightweight,rose2022ideres}.

Uma das formas mais diretas de contenção é o bloqueio e filtragem de tráfego, através de \textit{firewalls} ou listas de controlo de acessos (ACL). Estas permitem restringir comunicações com base em características como o endereço IP, a porta ou o protocolo utilizado. A resposta pode passar pelo descarte de pacotes, pelo bloqueio de determinados fluxos ou pela limitação da taxa de transmissão (\textit{rate limiting}). Uma intervenção mais abrangente consiste no isolamento do dispositivo, impedindo que este continue a comunicar livremente com os restantes equipamentos. Este isolamento pode ser conseguido através de regras de \textit{firewall}, listas negras (\textit{blacklists}) ou segmentação da rede, incluindo a utilização de VLANs de quarentena. A escolha entre uma restrição parcial e o isolamento completo deve considerar a gravidade da ameaça e o impacto sobre as comunicações legítimas~\cite{rose2022ideres,karibandi2025cyberdome,hunde2026rankshield}.

A contenção pode também ser realizada através de redes definidas por \textit{software} (SDN), que permitem alterar dinamicamente as políticas de encaminhamento e filtragem. Ao separar logicamente o plano de controlo do plano de dados, o SDN possibilita que um controlador central instale regras nos equipamentos de rede, recorrendo, por exemplo, ao OpenFlow. A utilização de \textit{switches} programáveis com P4 acrescenta a possibilidade de definir o processamento dos pacotes diretamente no plano de dados. Estas tecnologias permitem bloquear, redirecionar ou limitar tráfego malicioso sem exigir alterações ao \textit{software} dos dispositivos IoT~\cite{elsayed2024mpguard}.

O local onde a resposta é aplicada constitui outra decisão importante. A proteção pode funcionar no próprio dispositivo (\textit{host-based}), permitindo-lhe descartar pacotes suspeitos, embora esta abordagem esteja limitada pela capacidade de processamento e memória disponível~\cite{kponyo2020lightweight}. Em alternativa, a contenção pode ser executada num \textit{gateway} ou nó \textit{edge}, que concentra as comunicações de vários dispositivos e permite protegê-los sem instalar mecanismos de segurança individuais. Os controladores SDN e \textit{switches} programáveis permitem aplicar políticas em diferentes pontos da infraestrutura, enquanto a \textit{cloud} pode centralizar a análise e a tomada de decisões. Importa, por isso, distinguir o local onde a resposta é decidida daquele onde é efetivamente executada~\cite{elsayed2024mpguard,karibandi2025cyberdome}.

Por último, a eficácia da contenção depende também de como e quando a resposta é desencadeada. A estratégia mais simples consiste em aplicar uma medida predefinida após a classificação de tráfego como malicioso. Contudo, uma decisão baseada numa única observação pode originar bloqueios indevidos. Algumas abordagens recorrem, por isso, a limiares ou verificações adicionais antes de isolar um dispositivo~\cite{hunde2026rankshield}. Outras utilizam mecanismos de decisão, incluindo aprendizagem por reforço, para escolher ações de acordo com o estado da rede e as consequências esperadas~\cite{zolotukhin2021curious,rose2022ideres}. O momento da intervenção também é relevante, pois agir demasiado tarde pode permitir que o ataque produza danos, enquanto agir com informação insuficiente aumenta o risco de falsos positivos~\cite{tan2023hitting}. Assim, um mecanismo de contenção automática deve procurar equilibrar rapidez, fiabilidade e continuidade dos serviços legítimos.

% =====================================================================
\section{Metodologia da revisão sistemática}\label{sec:metodologia}
% RESPONSÁVEL: Frente 1 (Afonso)
% =====================================================================
A revisão segue o protocolo de RSL definido na unidade curricular. O processo teve
quatro fases: (1) definição das perguntas de pesquisa; (2) construção da \textit{string}
de busca e pesquisa nas bases de dados; (3) seleção dos estudos, com critérios de
inclusão e exclusão, leitura pelo método \textit{three-pass}~\cite{keshav2007read} e
avaliação de qualidade; e (4) extração dos dados para um formulário comum. Todo o
processo, da importação das referências à extração, foi gerido na ferramenta
Parsifal. Os ficheiros exportados de cada base, o registo das pesquisas
e os dados extraídos estão disponíveis no repositório do projeto~\cite{repo}.

\subsection{Perguntas de pesquisa}
Para delimitar o âmbito foi usada a estrutura PIPOC. A população são as redes e os
dispositivos IoT. A intervenção é a deteção de dispositivos comprometidos com ML a partir
do tráfego de rede. A comparação faz-se entre famílias de modelos e entre mecanismos de
resposta. Os resultados de interesse são o desempenho da deteção e a eficácia da
contenção. O contexto é a aplicação no \textit{gateway} ou na \textit{edge}, onde o
tráfego pode ser observado e bloqueado.

Daqui resultaram uma questão geral (QG), quatro questões focadas (QF1--QF4) e uma
questão estatística (QE), apresentadas na Tabela~\ref{tab:perguntas}. A QF1 e a QF2
caracterizam a deteção: que modelos se usam, com que dados são treinados e como são
avaliados. A QF3 é o centro da revisão, porque trata da resposta, isto é, do que o
sistema faz depois de detetar um dispositivo comprometido. A QF4 avalia se as propostas
são exequíveis fora de um \textit{dataset} estático. A QE dá a visão quantitativa do
conjunto de estudos.

\begin{table}[H]
\caption{Perguntas de pesquisa}\label{tab:perguntas}
\centering\footnotesize
\begin{tabular}{@{}p{0.09\columnwidth}p{0.83\columnwidth}@{}}
\toprule
ID & Pergunta \\
\midrule
QG  & Como é que a análise de tráfego com ML é usada para detetar e conter automaticamente dispositivos IoT comprometidos? \\
QF1 & Que modelos (supervisionados, não supervisionados e semi-supervisionados) e que \textit{features} de tráfego são usados? \\
QF2 & Que \textit{datasets}, tipos de ataque e métricas de avaliação são usados? \\
QF3 & Que mecanismos de resposta automática são propostos, onde são aplicados (dispositivo, \textit{gateway}/\textit{edge}, controlador SDN ou \textit{cloud}) e como são desencadeados? \\
QF4 & As soluções são avaliadas apenas \textit{offline}, com \textit{datasets}, ou também em tempo real ou num protótipo, e com que custo computacional? \\
QE  & Quantas publicações existem por ano, por tipo de modelo e por \textit{dataset}? \\
\bottomrule
\end{tabular}
\end{table}

As perguntas não ficaram fechadas à partida. Antes da triagem, a \textit{string} principal
foi complementada com os termos ``systematic literature'', ``systematic review'' e
``literature review'', para encontrar revisões já publicadas sobre o tema. Esta pesquisa
secundária devolveu dez resultados (oito no Scopus e dois no IEEE Xplore), que
correspondem a oito revisões distintas. A leitura das três mais próximas do
tema~\cite{kandah2025navigating,bhardwaj2025systematic,veramundaca2024cybersecurity},
complementada por uma revisão anterior sobre a deteção de \textit{botnets} em redes IoT
com SDN~\cite{negera2022review}, levou a três ajustes: a QF2 passou a incluir os tipos de ataque; a QF3 passou a perguntar
onde a resposta é aplicada e como é desencadeada; e foi acrescentada a QF4, para separar
as propostas avaliadas apenas sobre \textit{datasets} das que foram testadas num
protótipo ou em tempo real. As revisões encontradas servem de enquadramento e não entram
no conjunto de estudos analisados.

\subsection{\textit{String} de busca e bases de dados}
A \textit{string} de busca combina com AND quatro grupos de termos, cada um com os seus
sinónimos ligados por OR:
\begin{quote}\footnotesize\raggedright
(``internet of things'' OR ``IoT'') AND (``intrusion detection'' OR ``anomaly detection'' OR
``malware detection'' OR ``botnet detection'') AND (``machine learning'' OR ``deep learning'')
AND (``mitigation'' OR ``quarantine'' OR ``automated response'' OR ``intrusion prevention'' OR ``countermeasure'')
\end{quote}
O primeiro grupo fixa o domínio, o segundo a tarefa de deteção e o terceiro a técnica. O
quarto grupo, com os termos ligados à resposta, é o que distingue esta revisão de uma
revisão sobre deteção: sem ele, a pesquisa devolve sobretudo estudos que terminam no
alerta. Uma primeira versão da \textit{string}, que incluía também termos como
``network traffic'', ``traffic analysis'' e ``isolation'' e era aplicada a todos os
campos, devolvia 3184 resultados na ACM e 1910 no IEEE Xplore, na sua maioria fora do
tema. A versão final retira os termos mais genéricos e limita a pesquisa ao resumo. A
exigência de que a deteção use tráfego de rede passou para os critérios de exclusão
(EC3).

A pesquisa foi feita a 6 de outubro de 2026 em cinco bases (Tabela~\ref{tab:bases}). O
IEEE Xplore e a ACM Digital Library cobrem as principais conferências e revistas de redes
e segurança; o Scopus é um índice multidisciplinar e alarga a cobertura; o ScienceDirect
inclui as revistas da Elsevier, como a \textit{Computer Networks}; e a MDPI publica em
acesso aberto revistas como a \textit{Sensors}. Em todas foi pesquisado apenas o resumo,
para publicações entre 2020 e 2026. No Scopus foram mantidos apenas artigos de revista e
de conferência; na ACM foi usada a coleção \textit{ACM Guide to Computing Literature}. O
ScienceDirect limita o número de operadores booleanos por pesquisa, pelo que foi usada
uma versão reduzida, com os termos mais frequentes de cada grupo:
\begin{quote}\footnotesize\raggedright
``IoT'' AND (``intrusion detection'' OR ``anomaly detection'' OR ``botnet detection'')
AND (``machine learning'' OR ``deep learning'') AND (``mitigation'' OR
``intrusion prevention'' OR ``automated response'')
\end{quote}

\begin{table}[H]
\caption{Bases de dados pesquisadas (6/10/2026)}\label{tab:bases}
\centering\footnotesize
\begin{tabular}{@{}p{0.26\columnwidth}p{0.46\columnwidth}r@{}}
\toprule
Base & Filtros & Resultados \\
\midrule
IEEE Xplore         & resumo, 2020--2026 & 139 \\
ACM Digital Library & resumo, 2020--2026, ACM Guide & 34 \\
Scopus              & resumo, 2020--2026, \textit{article} e \textit{conference paper} & 342 \\
ScienceDirect       & resumo, 2020--2026, \textit{string} reduzida & 148 \\
MDPI                & resumo, 2020--2026 & 14 \\
\midrule
\textbf{Total}      & & \textbf{677} \\
\bottomrule
\end{tabular}
\end{table}

\subsection{Critérios de inclusão e exclusão}
Os critérios usados na seleção estão na Tabela~\ref{tab:criterios}. Os critérios de
inclusão garantem a qualidade mínima e a atualidade das fontes: estudos revistos por
pares, em inglês, publicados entre 2020 e 2026. Os critérios de exclusão afastam os
estudos que não servem as perguntas de pesquisa. O EC3 exclui as abordagens que dependem
de dados recolhidos no próprio dispositivo, como o consumo energético, chamadas ao
sistema ou análise de executáveis, porque não podem ser aplicadas num \textit{gateway}
que só vê o tráfego. O EC4 afasta as revisões e as propostas puramente conceptuais.

O EC6 não fazia parte do protocolo inicial. Foi acrescentado depois da primeira leitura
dos resumos, quando se percebeu que muitos estudos usavam termos como
\textit{mitigation} ou \textit{prevention} mas tratavam apenas a deteção, deixando a
resposta como trabalho futuro ou referindo-a numa frase. Como a QG pergunta como se
deteta \emph{e} se contém, esses estudos não respondem ao objetivo da revisão. O
refinamento está previsto no protocolo e ficou registado no repositório.

\begin{table}[H]
\caption{Critérios de inclusão (IC) e de exclusão (EC)}\label{tab:criterios}
\centering\footnotesize
\begin{tabular}{@{}p{0.09\columnwidth}p{0.83\columnwidth}@{}}
\toprule
ID & Critério \\
\midrule
IC1 & Publicado em revista, conferência ou \textit{workshop} com revisão por pares \\
IC2 & Escrito em inglês \\
IC3 & Publicado entre 2020 e 2026 \\
IC4 & Contém os termos da \textit{string} de busca \\
\midrule
EC1 & Não responde a nenhuma pergunta de pesquisa \\
EC2 & Não é sobre IoT ou ambientes semelhantes \\
EC3 & A deteção não se baseia em tráfego de rede \\
EC4 & Sem avaliação experimental \\
EC5 & Duplicado \\
EC6 & Não propõe nem avalia um mecanismo de resposta ou contenção automática \\
\bottomrule
\end{tabular}
\end{table}

\subsection{Processo de seleção}
Das 677 referências encontradas, 660 foram importadas para o Parsifal. A diferença
corresponde a registos que as bases não deixaram exportar (cinco no ScienceDirect) ou que
a ferramenta não conseguiu importar. A remoção de duplicados (EC5) eliminou 169
registos, correspondentes a artigos indexados em mais do que uma base, como o Scopus e o IEEE Xplore,
deixando 491 estudos únicos. Estes foram divididos em blocos pelos cinco elementos do
grupo, e cada bloco passou pelas etapas seguintes:
\begin{enumerate}
  \item \textbf{Título e resumo.} Aplicação dos critérios IC1--IC4 e EC1--EC4. Os casos
  duvidosos passaram à etapa seguinte, para não excluir estudos relevantes só pelo
  resumo.
  \item \textbf{Resposta automática.} Aplicação do EC6 aos estudos aceites. Quando o
  resumo não era conclusivo, a decisão foi tomada com o texto completo.
  \item \textbf{Primeira passagem e avaliação de qualidade.} Leitura de cada artigo pela
  primeira passagem do método \textit{three-pass}~\cite{keshav2007read}: título, resumo,
  introdução, títulos das secções e conclusões, complementada pelas secções de método e
  de resultados necessárias para responder às perguntas de qualidade da
  Tabela~\ref{tab:qualidade}.
\end{enumerate}

Cada pergunta de qualidade tem três respostas possíveis, com o valor
$q_k \in \{0;\ 0{,}5;\ 1\}$ (Não, Parcialmente, Sim). A pontuação de um estudo $e$ é a soma
\begin{equation}
  Q(e) = \sum_{k=1}^{8} q_k(e), \qquad 0 \le Q(e) \le 8,
  \label{eq:qualidade}
\end{equation}
e só foram incluídos os estudos com $Q(e) \ge 5$. As perguntas QA2 a QA7 estão ligadas às
perguntas de pesquisa (Tabela~\ref{tab:qualidade}), o que garante que os estudos
incluídos têm informação suficiente para lhes responder. A QA5 é deliberadamente exigente,
porque muitos artigos anunciam um mecanismo de resposta mas não o implementam nem
apresentam resultados sobre ele. A Fig.~\ref{fig:selecao} resume o número de estudos em
cada etapa. Dos 491 estudos únicos, 427 foram excluídos na seleção, pelos critérios
EC1--EC4 ou pelo EC6. Dos 64 que chegaram à avaliação de qualidade, 33 ficaram abaixo do
limiar e 31 foram incluídos.

\begin{table}[H]
\caption{Perguntas da avaliação de qualidade}\label{tab:qualidade}
\centering\footnotesize
\begin{tabular}{@{}p{0.09\columnwidth}p{0.68\columnwidth}p{0.12\columnwidth}@{}}
\toprule
ID & Pergunta & Relação \\
\midrule
QA1 & Os objetivos e a contribuição do estudo estão claramente definidos? & -- \\
QA2 & O modelo de ML e as \textit{features} de tráfego estão descritos com detalhe suficiente para se reproduzir? & QF1 \\
QA3 & O estudo usa um \textit{dataset} identificado e público, ou disponibiliza o seu? & QF2 \\
QA4 & Reporta métricas além da \textit{accuracy}, como \textit{precision}, \textit{recall}, F1 ou FPR? & QF2 \\
QA5 & O mecanismo de resposta ou contenção é descrito e avaliado, e não só mencionado? & QF3 \\
QA6 & A solução é testada em tempo real, num protótipo ou num \textit{testbed}, e não só \textit{offline}? & QF4 \\
QA7 & Reporta o custo computacional, como latência, CPU ou memória? & QF4 \\
QA8 & Discute as limitações do trabalho? & -- \\
\bottomrule
\end{tabular}
\end{table}

\begin{figure}[H]
\centering
\begin{tikzpicture}[
  font=\scriptsize, node distance=3.2mm,
  etapa/.style={draw, rounded corners=1pt, align=center, text width=0.47\columnwidth, inner sep=2pt},
  saida/.style={draw, dashed, align=left, text width=0.33\columnwidth, inner sep=2pt},
  seta/.style={-{Stealth[length=1.6mm]}, thin}]
  \node[etapa] (id)  {Identificação: IEEE 139, ACM 34, Scopus 342,\\ScienceDirect 148, MDPI 14 ($n=677$)};
  \node[etapa, below=of id]  (imp) {Importados no Parsifal ($n=660$)};
  \node[etapa, below=of imp] (uni) {Estudos únicos ($n=491$)};
  \node[etapa, below=of uni] (ec6) {Aceites na seleção: título, resumo e EC6 ($n=64$)};
  \node[etapa, below=of ec6, fill=gray!12] (inc) {Incluídos após avaliação de qualidade ($n=31$)};
  \draw[seta] (id) -- (imp);  \draw[seta] (imp) -- (uni); \draw[seta] (uni) -- (ec6);
  \draw[seta] (ec6) -- (inc);
  \node[saida, right=1.5mm of imp] (x0) {Não exportados ou não importados ($n=17$)};
  \node[saida, right=1.5mm of uni] (x1) {Duplicados, EC5 ($n=169$)};
  \node[saida, right=1.5mm of ec6] (x3) {Excluídos, EC1--EC4 e EC6 ($n=427$)};
  \node[saida, right=1.5mm of inc] (x4) {$Q(e)<5$ ($n=33$)};
\end{tikzpicture}
\caption{Processo de seleção dos estudos.}\label{fig:selecao}
\end{figure}

\subsection{Extração de dados}
Para cada estudo incluído foi preenchido no Parsifal um formulário comum, com campos de
texto livre e campos de escolha fixa. Os campos de escolha fixa (tipo de aprendizagem,
local da deteção, mecanismo de resposta, tipos de ataque e tipo de avaliação) usam as
mesmas categorias para todos os estudos, o que permite contá-los e compará-los
diretamente. A Tabela~\ref{tab:extracao} mostra como os campos se distribuem pelas
perguntas de pesquisa. Os dados extraídos foram exportados do Parsifal para uma folha de
cálculo, disponível no repositório~\cite{repo}, a partir do qual foram construídos os gráficos e
as tabelas da Secção~\ref{sec:resultados}.

\begin{table}[H]
\caption{Campos do formulário de extração}\label{tab:extracao}
\centering\footnotesize
\begin{tabular}{@{}p{0.72\columnwidth}p{0.18\columnwidth}@{}}
\toprule
Campos & Pergunta \\
\midrule
Referência BibTeX, DOI, ano, autores, tipo e local de publicação & QE \\
Modelos de ML/DL, tipo de aprendizagem, \textit{features} de tráfego & QF1 \\
\textit{Datasets}, tipos de ataque detetados, métricas e resultados & QF2 \\
Local da deteção, mecanismo de resposta, forma de desencadear a resposta & QF3 \\
Tipo de avaliação, tempo real, \textit{hardware} de teste, latência e custo computacional & QF4 \\
Principais limitações & Discussão \\
\bottomrule
\end{tabular}
\end{table}

% =====================================================================
\section{Resultados da revisão}\label{sec:resultados}
% Cada subsecção responde a UMA pergunta de pesquisa.
% =====================================================================
\subsection{Visão geral dos estudos (QE)}
% RESPONSÁVEL: Frente 1 (Afonso)
O conjunto final é composto por 31 estudos, todos
com pontuação de qualidade $Q(e) \ge 5$ (Eq.~\ref{eq:qualidade}). Vinte e quatro foram publicados em revista e sete em conferência. A distribuição por ano
(Fig.~\ref{fig:anos}) mostra que o tema é recente e está a crescer: 24 dos 31 estudos são
de 2024 ou posterior.

\begin{figure}[H]
\centering
\begin{tikzpicture}
\begin{axis}[
  ybar, bar width=9pt, width=0.95\columnwidth, height=3.6cm,
  ymin=0, ymax=14, ytick={0,3,6,9,12},
  symbolic x coords={2020,2021,2022,2023,2024,2025,2026}, xtick=data,
  nodes near coords, nodes near coords style={font=\scriptsize},
  tick label style={font=\scriptsize}, ylabel={\scriptsize Estudos},
  axis lines*=left, enlarge x limits=0.08, ymajorgrids, grid style={gray!25}]
\addplot[fill=gray!55, draw=black] coordinates
  {(2020,2) (2021,2) (2022,3) (2023,0) (2024,6) (2025,10) (2026,8)};
\end{axis}
\end{tikzpicture}
\caption{Estudos incluídos por ano de publicação.}\label{fig:anos}
\end{figure}

Quanto ao tipo de modelo (Fig.~\ref{fig:tipos}), 29 dos 31 estudos usam aprendizagem
supervisionada, cinco recorrem também a aprendizagem por reforço, em geral para escolher
a resposta, três usam aprendizagem semi-supervisionada e só dois incluem modelos não
supervisionados. Quanto aos \textit{datasets} (Fig.~\ref{fig:datasets_qe}), o UNSW-NB15 e o
IoT-23 são os mais usados, em cinco estudos cada, seguidos do ToN-IoT, do Edge-IIoTset e
do CICIoT2023, com quatro cada, e 17 estudos recorrem também a dados próprios. As Secções~\ref{sec:qf1} e~\ref{sec:qf2}
analisam estes resultados em detalhe.

\begin{figure}[H]
\centering
\begin{tikzpicture}
\begin{axis}[
  xbar, bar width=6pt, width=0.78\columnwidth, height=3.0cm,
  xmin=0, xmax=35, xtick={0,10,20,30},
  symbolic y coords={Não supervisionada,Semi-supervisionada,Por reforço,Supervisionada}, ytick=data,
  nodes near coords, nodes near coords style={font=\scriptsize},
  tick label style={font=\scriptsize}, xlabel={\scriptsize Estudos},
  axis lines*=left, enlarge y limits=0.2, xmajorgrids, grid style={gray!25}]
\addplot[fill=gray!55, draw=black] coordinates
  {(29,Supervisionada) (5,Por reforço) (3,Semi-supervisionada) (2,Não supervisionada)};
\end{axis}
\end{tikzpicture}
\caption{Estudos por tipo de aprendizagem (um estudo pode usar mais do que um).}\label{fig:tipos}
\end{figure}

\begin{figure}[H]
\centering
\begin{tikzpicture}
\begin{axis}[
  xbar, bar width=5pt, width=0.78\columnwidth, height=4.6cm,
  xmin=0, xmax=22, xtick={0,5,10,15,20},
  symbolic y coords={X-IIoTID,CIC-IDS2017,CICIDS2018,BoT-IoT,CICIoT2023,Edge-IIoTset,ToN-IoT,IoT-23,UNSW-NB15,Dados próprios}, ytick=data,
  nodes near coords, nodes near coords style={font=\scriptsize},
  tick label style={font=\scriptsize}, xlabel={\scriptsize Estudos},
  axis lines*=left, enlarge y limits=0.07, xmajorgrids, grid style={gray!25}]
\addplot[fill=gray!55, draw=black] coordinates
  {(17,Dados próprios) (5,UNSW-NB15) (5,IoT-23) (4,ToN-IoT) (4,Edge-IIoTset) (4,CICIoT2023) (3,BoT-IoT) (3,CICIDS2018) (3,CIC-IDS2017) (3,X-IIoTID)};
\end{axis}
\end{tikzpicture}
\caption{\textit{Datasets} mais usados pelos estudos incluídos.}\label{fig:datasets_qe}
\end{figure}

A Tabela~\ref{tab:qa_resultados} mostra as respostas de cada estudo às oito perguntas de
qualidade, e a última linha conta quantos estudos respondem ``sim'' a cada uma. As
pontuações vão de 5,0 a 7,5. Os objetivos estão claros em todos os estudos (QA1) e quase
todos reportam métricas além da \textit{accuracy} (QA4). As perguntas mais fracas são,
precisamente, as que mais interessam a esta revisão: só 14 dos 31 estudos descrevem
\emph{e} avaliam o mecanismo de resposta (QA5), 17 testam a solução em tempo real ou num
protótipo (QA6) e 16 reportam o custo computacional (QA7). Também a reprodutibilidade
(QA2) e a discussão das limitações (QA8) ficam quase sempre pela metade. Durante a
leitura foram ainda encontradas incoerências numéricas internas em vários
estudos~\cite{elsayed2024mpguard,ibrahim2026purechain,andrade2024ddos,dange2025dual,hussain2026postquantum,premakumari2025qlearning,kumar2025advanced},
como valores que mudam entre secções ou métricas que não batem com a matriz de confusão
apresentada. Por isso, os valores reportados são usados nesta revisão tal como os
autores os publicaram, mas comparados com cautela.

\begin{table}[H]
\caption{Avaliação de qualidade dos estudos incluídos (S: sim; P: parcialmente; N: não)}\label{tab:qa_resultados}
\centering\scriptsize
\setlength{\tabcolsep}{3.2pt}
\begin{tabular}{@{}lcccccccc c@{}}
\toprule
Estudo & QA1 & QA2 & QA3 & QA4 & QA5 & QA6 & QA7 & QA8 & $Q(e)$ \\
\midrule
Al Abdulwahid~\cite{Abdulwahid2025} & S & S & S & S & S & S & P & S & 7,5 \\
Karam~\cite{karam2025isaaf} & S & S & S & S & S & S & S & P & 7,5 \\
Toony~\cite{toony2024multiblock} & S & P & S & S & S & S & S & S & 7,5 \\
Arulkumar~\cite{Arulkumar2026}      & S & S & S & S & P & S & S & P & 7,0 \\
Hussain~\cite{hussain2026postquantum} & S & S & S & S & P & P & S & S & 7,0 \\
Iqbal~\cite{Iqbal2026}              & S & P & S & S & P & S & S & S & 7,0 \\
Kponyo~\cite{kponyo2020lightweight}   & S & S & S & N & S & S & S & S & 7,0 \\
Ramirez-Martinez~\cite{ramirezmartinez2025p4} & S & P & S & S & S & S & S & P & 7,0 \\
Shukla~\cite{shukla2026cyber} & S & S & S & S & P & P & S & S & 7,0 \\
Srilekha~\cite{srilekha2025hybrid} & S & P & S & S & S & S & S & P & 7,0 \\
Andrade~\cite{andrade2024ddos}        & S & P & S & S & S & S & S & N & 6,5 \\
Gowtham~\cite{gowtham2026fedmp} & S & P & P & S & S & S & S & P & 6,5 \\
Mazhar~\cite{mazhar2022ridps} & S & P & P & S & S & S & S & P & 6,5 \\
Ravi~\cite{ravi2020semisupervised}    & S & P & P & S & S & S & S & P & 6,5 \\
Rose~\cite{rose2022ideres}            & S & P & S & S & S & S & P & P & 6,5 \\
Chaudhary~\cite{Chaudhary2025}      & S & S & S & S & P & P & P & P & 6,0 \\
Garba~\cite{garba2024sdn} & S & P & S & S & P & S & P & P & 6,0 \\
Ibrahim~\cite{ibrahim2026purechain}   & S & P & S & S & P & P & S & P & 6,0 \\
Premakumari~\cite{premakumari2025qlearning} & S & P & P & P & S & S & S & P & 6,0 \\
Rasheed~\cite{Rasheed2026}          & S & P & S & S & P & N & S & S & 6,0 \\
Saurabh~\cite{saurabh2022mlaps}       & S & S & S & S & P & N & P & S & 6,0 \\
Tawfik~\cite{tawfik2021internet}      & S & P & S & S & P & S & N & S & 6,0 \\
Verma~\cite{verma2024petrak} & S & S & S & S & P & P & P & P & 6,0 \\
El-Sayed~\cite{elsayed2024mpguard}    & S & P & S & S & P & P & P & P & 5,5 \\
Hunde~\cite{hunde2026rankshield}      & S & P & P & S & S & P & P & P & 5,5 \\
Tariq~\cite{tariq2025saegru} & S & P & S & S & P & P & P & P & 5,5 \\
Dange~\cite{dange2025dual}            & S & P & S & S & N & P & P & P & 5,0 \\
Karibandi~\cite{karibandi2025cyberdome} & S & P & N & S & P & S & P & P & 5,0 \\
Kumar e Kumar~\cite{kumar2025advanced} & S & P & S & S & P & N & P & P & 5,0 \\
Rawat~\cite{rawat2024intelligent}     & S & S & S & S & N & N & N & S & 5,0 \\
Zolotukhin~\cite{zolotukhin2021curious} & S & P & S & P & S & P & N & P & 5,0 \\
\midrule
N.º de ``S'' (31) & 31 & 10 & 25 & 28 & 14 & 17 & 16 & 10 & \\
\bottomrule
\end{tabular}
\end{table}

\subsection{Modelos e \textit{features} (QF1)}\label{sec:qf1}
% RESPONSÁVEL: Frente 3 (Rafael Silva, supervisionado e features) + Frente 4 (Rafael Reis, não supervisionado)
% Parte supervisionada da QF1 (Rafael Silva; integrada com os restantes estudos).
A abordagem supervisionada é a mais comum nos estudos analisados. Os modelos clássicos
continuam presentes, sobretudo a RF, as árvores de decisão, o KNN, a regressão logística
(LR), o Naive Bayes e a \textit{Support Vector Machine}
(SVM)~\cite{ramirezmartinez2025p4,andrade2024ddos,dange2025dual,kponyo2020lightweight,saurabh2022mlaps,rawat2024intelligent,karam2025isaaf,hussain2026postquantum,mazhar2022ridps,garba2024sdn,verma2024petrak},
ao lado de redes profundas como CNN, LSTM, GRU e
CNN-BiLSTM~\cite{rose2022ideres,ibrahim2026purechain,tawfik2021internet,srilekha2025hybrid}.

Muitos estudos utilizam vários modelos, cada um de forma diferente: o
\textit{stacking} de seis classificadores com um MLP como
meta-modelo~\cite{elsayed2024mpguard} ou um \textit{ensemble} ponderado de seis
classificadores~\cite{toony2024multiblock}; a combinação de SVM, RF, J48,
\textit{K-Nearest Neighbor} (KNN) e \textit{Extreme Gradient Boosting} (XGBoost) num único
sistema de deteção de ataques~\cite{Chaudhary2025}; a junção de XGBoost com LR num
meta-classificador, num sistema de deteção multinível que integra um \textit{stacked
ensemble}~\cite{Arulkumar2026}; uma rede neuronal de grafos (GNN) combinada com LSTM para analisar as relações espaciais
e temporais de um sistema ciber-físico (CPS)~\cite{Rasheed2026}, ou um \textit{graph
transformer} com o mesmo fim~\cite{shukla2026cyber}; um classificador baseado em RF para
detetar os ataques, seguido de um classificador MLP especializado que identifica o tipo
de ataque~\cite{Iqbal2026}; Extra Trees no \textit{edge} e CNN-LSTM na
\textit{cloud}~\cite{srilekha2025hybrid}; uma Bi-GRU com atenção~\cite{kumar2025advanced}; uma
rede Ghost-SE treinada com aprendizagem federada~\cite{gowtham2026fedmp}; XGBoost~\cite{karibandi2025cyberdome};
e uma rede neuronal profunda combinada com limiares
dinâmicos~\cite{hunde2026rankshield}. Cinco estudos usam ainda
aprendizagem por reforço para escolher a
resposta~\cite{zolotukhin2021curious,gowtham2026fedmp,premakumari2025qlearning,kumar2025advanced,Arulkumar2026}.

As \textit{features} de tráfego são um conceito importante para a deteção de anomalias e
de ataques em todos estes estudos e, por esta razão, os artigos apresentam as
\textit{features} que o modelo considerou importantes durante a sua fase de teste.
Al Abdulwahid, por exemplo, destaca o tamanho dos
pacotes (\textit{payload size}), o tempo entre chegadas de pacotes
(\textit{inter-arrival time}), a duração do fluxo (\textit{flow duration}), o número de
tentativas de ligação falhadas, a entropia dos IP de origem e de destino e o protocolo
utilizado~\cite{Abdulwahid2025}; Arulkumar e K, a frequência de comandos PLC, os picos de
transmissão de dados e as anomalias de autenticação~\cite{Arulkumar2026}. De forma geral, a maioria dos estudos trabalha com
estatísticas de fluxo, como contagens de pacotes e de \textit{bytes}, duração e intervalos
entre pacotes, calculadas por fluxo ou por janelas de tempo curtas, de 0,5\,s a
2\,s~\cite{ravi2020semisupervised,elsayed2024mpguard}, e três estudos calculam as
\textit{features} no próprio plano de dados, em \textit{switches}
P4~\cite{ramirezmartinez2025p4,andrade2024ddos,elsayed2024mpguard}.


% Parte não supervisionada da QF1 (Rafael Reis).
As abordagens que não dependem só de rótulos são minoritárias. Ravi e
Shalinie~\cite{ravi2020semisupervised} combinam uma rede neuronal com K-means num esquema
semi-supervisionado que treina com apenas 10\% dos dados rotulados e marca como incertos
os casos em que os modelos discordam. Zolotukhin \textit{et al.}~\cite{zolotukhin2021curious}
usam \textit{autoencoders}, mapas auto-organizados e BiGAN para a deteção e um agente de
aprendizagem por reforço para decidir a resposta. Tariq e Ahanger treinam um
\textit{autoencoder} empilhado sem rótulos antes de um classificador
GRU~\cite{tariq2025saegru}, Shukla e Agarwal usam pré-treino
auto-supervisionado~\cite{shukla2026cyber}, e Al Abdulwahid opta por uma aprendizagem
semi-supervisionada, que combina \textit{autoencoders} para a deteção de anomalias com
redes neuronais recorrentes (RNN) para a análise temporal, detetando anomalias tanto
imediatas como de longo prazo~\cite{Abdulwahid2025}.
Mesmo nos estudos que juntam modelos supervisionados e não
supervisionados~\cite{zolotukhin2021curious,tariq2025saegru,Abdulwahid2025}, os dois são
usados em conjunto ou a escolha entre eles fica a cargo de um agente de reforço, sem uma
análise do efeito de cada um na contenção. Essa comparação é um dos
objetivos da solução proposta na Secção~\ref{sec:proposta}.

As aplicações destes modelos distinguem-se pelos recursos de \textit{hardware} e pela
complexidade do tráfego analisado. Os modelos leves, como o \textit{Isolation Forest} ou
heurísticas com aprendizagem automática simples, são os que cabem nos próprios
dispositivos ou em \textit{gateways} de baixo custo. É o caso do mecanismo de Kponyo
\textit{et al.}~\cite{kponyo2020lightweight}, que deteta \textit{floods} SYN, ICMP e UDP
diretamente em dispositivos IoT com poucos recursos. O modelo, um KNN treinado fora do
dispositivo, classifica o efeito de um afluxo de pacotes na CPU e na memória e ocupa em
média 23\,MB. Os modelos profundos, como os \textit{autoencoders} ou as redes CNN e LSTM,
exigem mais capacidade de treino e aparecem sobretudo em arquiteturas centralizadas, em
controladores SDN ou em servidores na \textit{cloud}, onde analisam os dados agregados de
vários
dispositivos~\cite{zolotukhin2021curious,tawfik2021internet,rose2022ideres}.
Esta complementaridade entre modelos leves para a reação imediata e modelos profundos para
a análise centralizada é uma das razões para a solução proposta comparar o
\textit{Isolation Forest} e o \textit{autoencoder} no mesmo \textit{gateway}.

% Síntese da QF1 (Rafael Silva).
Em síntese, a aprendizagem supervisionada é a mais comum entre os estudos analisados;
alguns poucos seguem a semi-supervisionada, menos ainda a não supervisionada, e alguns
aproveitam ideias da aprendizagem por reforço nos seus
modelos~\cite{Arulkumar2026,zolotukhin2021curious,gowtham2026fedmp,premakumari2025qlearning,kumar2025advanced}. Da mesma forma, há uma grande
diferença entre as \textit{features} de tráfego consideradas importantes pelos vários
modelos, e diversos estudos não as indicam, ou explicam apenas como o modelo decidiu
quais eram as mais importantes, sem nunca as
listar~\cite{Rasheed2026,Chaudhary2025}. Das que foram indicadas, algumas das
mais comuns são o tamanho dos pacotes, que se altera quando o dispositivo envia código
malicioso ou comandos fora do padrão; o tempo entre chegadas dos pacotes, que num ataque
DDoS se torna muito curto, por causa do volume de pacotes enviados; e a duração do fluxo,
que pode indicar comunicação maliciosa se for muito diferente da habitual entre os
dispositivos.

\subsection{\textit{Datasets}, tipos de ataque e métricas (QF2)}\label{sec:qf2}
% RESPONSÁVEL: Frente 2 (João Alves)
\begin{table*}[p]
\caption{Datasets usados por estudo}\label{tab:datasets}
\centering\scriptsize
\renewcommand{\arraystretch}{0.84}
\setlength{\tabcolsep}{3pt}
\begin{tabular}{@{}>{\raggedright\arraybackslash}p{0.19\textwidth}>{\raggedright\arraybackslash}p{0.36\textwidth}>{\raggedright\arraybackslash}p{0.41\textwidth}@{}}
\toprule
Estudo & \textit{Datasets} públicos & Dados próprios ou simulados \\
\midrule
Al Abdulwahid~\cite{Abdulwahid2025} & CICIDS2017, CSE-CIC-IDS2018, UNSW-NB15 & Nenhum \\
Andrade~\cite{andrade2024ddos} & CICDDoS2019 & Nenhum \\
Arulkumar~\cite{Arulkumar2026} & DAPT2020, IoT Intrusion Detection, KDDCUP'99, CICIoT2023 & Nenhum \\
Chaudhary~\cite{Chaudhary2025} & BoT-IoT & Nenhum \\
Dange~\cite{dange2025dual} & BoT-IoT, CTU-13 & \textit{Testbed} próprio, a pedido \\
El-Sayed~\cite{elsayed2024mpguard} & ToN-IoT, Edge-IIoTset, CICIoT2023 & Nenhum \\
Garba~\cite{garba2024sdn} & UNSW-NB15, CICDDoS2019 & Captura de 24\,h em \textit{testbed}, a pedido \\
Gowtham (FedMP-LDL)~\cite{gowtham2026fedmp} & Edge-IIoTset, CIC-IDS2017 (só validação cruzada) & \textit{Testbed} Mininet SDN-\textit{fog}, não disponibilizado \\
Hunde (RankShield)~\cite{hunde2026rankshield} & -- & Simulação em Cooja com ataques ao RPL, 70/30 \\
Hussain~\cite{hussain2026postquantum} & NSL-KDD, MedBIoT, CIC-IDS2017, Edge-IIoTset, IoT-23, ToN-IoT & Nenhum \\
Ibrahim~\cite{ibrahim2026purechain} & IoTForge Pro, WUSTL-IIoT-2021, X-IIoTID & Nenhum (reprodução dos \textit{datasets}) \\
Iqbal (DeepGuard)~\cite{Iqbal2026} & WUSTL-IIoT-2021, X-IIoTID & Nenhum \\
Karam (ISAAF)~\cite{karam2025isaaf} & MQTTset & MQTTEEB-D: \textit{testbed} real de IoT de saúde, 222\,813 registos, público no Mendeley \\
Karibandi (CyberDome)~\cite{karibandi2025cyberdome} & -- & \textit{Dataset} com 11 classes, de origem não identificada \\
Kponyo~\cite{kponyo2020lightweight} & -- & Medições de CPU e memória no \textit{testbed} (30\,384 amostras), não públicas \\
Kumar e Kumar~\cite{kumar2025advanced} & CICIoT2023 & Nenhum \\
Mazhar (R-IDPS)~\cite{mazhar2022ridps} & -- & \textit{Testbed} Mininet-WiFi com MQTT e ONOS, não público \\
Premakumari~\cite{premakumari2025qlearning} & AWID, IoT-23, WSN-BFSF (uso pouco claro) & Simulação em MATLAB Simulink com 50 nós \\
Ramirez-Martinez~\cite{ramirezmartinez2025p4} & CICIoT2023; IoT-23 (só tráfego benigno) & Tráfego Slowloris próprio, a pedido \\
Rasheed (XAI-CPS)~\cite{Rasheed2026} & SWaT, WADI, BATADAL, IIoT-Security & Nenhum \\
Ravi~\cite{ravi2020semisupervised} & NSL-KDD & \textit{Testbed} com 115\,200 registos (10\% rotulados), não público \\
Rawat~\cite{rawat2024intelligent} & IoT-23 & Nenhum \\
Rose~\cite{rose2022ideres} & -- & 913 PCAPs e imagens, públicos no IEEE DataPort \\
Saurabh~\cite{saurabh2022mlaps} & UNSW-NB15 & Nenhum \\
Shukla~\cite{shukla2026cyber} & UNSW-NB15, CICIDS2018, BoT-IoT & Smart-ITS, simulado com SUMO, origem pouco clara \\
Srilekha~\cite{srilekha2025hybrid} & UNSW-NB15 (aumentado com padrões IoT) & Telemetria IoT gerada com anomalias sintéticas \\
Tariq (SAE-GRU)~\cite{tariq2025saegru} & IoT-23, MedBIoT & Emulação de cidade inteligente (mais de 200 nós), parcialmente no Zenodo \\
Tawfik~\cite{tawfik2021internet} & IoTID20 & Tráfego emulado em Mininet, não público \\
Toony (MULTI-BLOCK)~\cite{toony2024multiblock} & Edge-IIoTset, ToN-IoT, X-IIoTID & Tráfego emulado em Mininet-WiFi para os testes de ataque, a pedido \\
Verma (PETRAK)~\cite{verma2024petrak} & KDDCUP'99, CIC-MalMem-2022, ToN-IoT, CICEV2023, VeReMi & Simulação em NS-3 com 697\,792 registos, a pedido \\
Zolotukhin~\cite{zolotukhin2021curious} & CICIDS2018 & Nenhum \\
\bottomrule
\end{tabular}

\vspace{3mm}
\caption{Tipos de ataque detetados por estudo (DoS: DoS/DDoS; Bot: botnet e C\&C; FB: força bruta; Mal: malware)}\label{tab:ataques}
\centering\scriptsize
\setlength{\tabcolsep}{4pt}\renewcommand{\arraystretch}{0.84}
\begin{tabular}{@{}>{\raggedright\arraybackslash}p{0.17\textwidth}cccc>{\raggedright\arraybackslash}p{0.62\textwidth}@{}}
\toprule
Estudo & DoS & Bot & FB & Mal & Outros ataques \\
\midrule
Al Abdulwahid~\cite{Abdulwahid2025} & x & x & x & x & \\
Andrade~\cite{andrade2024ddos} & x & & & & Não \\
Arulkumar~\cite{Arulkumar2026} & x & x & & & Ameaças persistentes avançadas (APT) \\
Chaudhary~\cite{Chaudhary2025} & x & & & & Não \\
Dange~\cite{dange2025dual} & & x & & & Varrimento \\
El-Sayed~\cite{elsayed2024mpguard} & x & & x & x & Varrimento, \textit{spoofing}, MITM, injeção \textit{web} \\
Garba~\cite{garba2024sdn} & x & & & & Não (SYN \textit{flood} com IP falsificado) \\
Gowtham~\cite{gowtham2026fedmp} & x & & & & Não \\
Hunde~\cite{hunde2026rankshield} & & & & & Ataques ao \textit{rank} do RPL \\
Hussain~\cite{hussain2026postquantum} & x & x & & x & MITM, \textit{spoofing}, reconhecimento \\
Ibrahim~\cite{ibrahim2026purechain} & x & & x & x & Varrimento, reconhecimento, XSS \\
Iqbal~\cite{Iqbal2026} & x & x & & x & \\
Karam~\cite{karam2025isaaf} & x & & x & & MQTT: Slowite, dados malformados, \textit{publish flood} \\
Karibandi~\cite{karibandi2025cyberdome} & x & x & x & x & Varrimento, outros \\
Kponyo~\cite{kponyo2020lightweight} & x & & & & Não (\textit{floods} ICMP, UDP, SYN) \\
Kumar e Kumar~\cite{kumar2025advanced} & x & & & x & SQLi, \textit{buffer overflow} \\
Mazhar~\cite{mazhar2022ridps} & x & & & & Não (\textit{floods} ICMP e SYN) \\
Premakumari~\cite{premakumari2025qlearning} & x & & & & \textit{Black-hole}, \textit{wormhole}, \textit{selective forwarding}, injeção de dados \\
Ramirez-Martinez~\cite{ramirezmartinez2025p4} & x & & & & Não (Slowloris) \\
Rasheed~\cite{Rasheed2026} & x & x & & & \\
Ravi~\cite{ravi2020semisupervised} & x & & & & Varrimento; R2L e U2R (NSL-KDD) \\
Rawat~\cite{rawat2024intelligent} & & x & & & Não \\
Rose~\cite{rose2022ideres} & x & x & x & x & \textit{Exploits} de serviços, enumeração \\
Saurabh~\cite{saurabh2022mlaps} & x & & & x & Reconhecimento, \textit{fuzzers} \\
Shukla~\cite{shukla2026cyber} & x & x & x & x & Reconhecimento, exfiltração, \textit{zero-day} \\
Srilekha~\cite{srilekha2025hybrid} & x & & & x & Varrimento, \textit{insider}, \textit{fuzzers} \\
Tariq~\cite{tariq2025saegru} & x & x & x & x & Varrimento, \textit{spoofing}, exfiltração, escuta \\
Tawfik~\cite{tawfik2021internet} & x & x & & & Varrimento, MITM (ARP \textit{spoofing}) \\
Toony~\cite{toony2024multiblock} & x & x & & & \textit{Spoofing}, MITM, varrimento, SQLi, XSS \\
Verma~\cite{verma2024petrak} & x & x & & x & \textit{Backdoor}, injeção, \textit{ransomware}, varrimento \\
Zolotukhin~\cite{zolotukhin2021curious} & x & x & x & & XSS, SQLi \\
\midrule
Total (31 estudos) & 28 & 15 & 9 & 13 & \\
\bottomrule
\end{tabular}
\end{table*}
\begin{table*}[p]
\caption{Métricas reportadas e melhores resultados por estudo (P: precision; R: recall; Acc: accuracy)}\label{tab:metricas}
\centering\scriptsize
\renewcommand{\arraystretch}{0.84}
\setlength{\tabcolsep}{3pt}
\begin{tabular}{@{}>{\raggedright\arraybackslash}p{0.13\textwidth}>{\raggedright\arraybackslash}p{0.30\textwidth}>{\raggedright\arraybackslash}p{0.53\textwidth}@{}}
\toprule
Estudo & Métricas reportadas & Melhor resultado reportado \\
\midrule
Al Abdulwahid~\cite{Abdulwahid2025} & Acc, P, R, F1, FPR & Acc 96,1\%, F1 95,0\%, FPR 3,2\% \\
Andrade~\cite{andrade2024ddos} & Acc, R, F1 & NB: Acc 98,04\%, F1 98,98\% (\textit{offline}) \\
Arulkumar~\cite{Arulkumar2026} & Acc, P, F1, especificidade, sensibilidade, NPV, MCC, FPR, FNR & Acc 97,67\% (divisão 70:30) e 98,92\% (outra divisão); FPR 0,0227 \\
Chaudhary~\cite{Chaudhary2025} & Acc, P, R, F1, FPR & RF: Acc 98,89\%, F1 97,74\%, FPR 0,006 \\
Dange~\cite{dange2025dual} & Acc, P, R, F1 & SVM: Acc 99,97\% no BoT-IoT; 0,92 no \textit{testbed} \\
El-Sayed~\cite{elsayed2024mpguard} & Acc, F1, FPR, AUC & CICIoT2023: Acc 99,32\%, F1 99,23\%, FPR 0,49\% \\
Garba~\cite{garba2024sdn} & Acc, P, R, F1 & DT: 99,57\% (próprio), 98,20\% (UNSW-NB15), 99,95\% (CICDDoS2019) \\
Gowtham~\cite{gowtham2026fedmp} & Acc, P, R, F1, AUC, FPR & Acc 99,72\% (\textit{dataset} não explicitado); 96,87\% no \textit{testbed}; FPR 0,0032\% (0,83\% no \textit{testbed}) \\
Hunde~\cite{hunde2026rankshield} & Acc, P, R, F1; latência, energia e PDR & Acc até 99,8\%, F1 até 99,7\%; PDR +67\%, latência $-$66\% \\
Hussain~\cite{hussain2026postquantum} & F1, Acc, P, R & RF: F1 0,9998--1,000 em cinco \textit{datasets}, mas 0,751 no NSL-KDD (média 0,957) \\
Ibrahim~\cite{ibrahim2026purechain} & Acc, F1, taxa de sucesso dos ataques & Acc 97,79--99,76\%; F1 99,1\% (multiclasse); sucesso dos ataques 0--8,79\% \\
Iqbal~\cite{Iqbal2026} & Acc, FPR & Acc 99,99\% (\textit{desktop}), 99,82\% (Raspberry Pi 5); FPR 0,2\% \\
Karam~\cite{karam2025isaaf} & Acc, P, R, F1, AUC & Acc 0,851 no MQTTset; cai para 0,08--0,21 noutro \textit{dataset}; 0,87 após novo treino \\
Karibandi~\cite{karibandi2025cyberdome} & Acc, F1 (macro) & Acc 99,99\%, F1 0,999 \\
Kponyo~\cite{kponyo2020lightweight} & Acc & KNN: 100\% \\
Kumar e Kumar~\cite{kumar2025advanced} & Acc, P, sensibilidade, especificidade, F, FPR, FNR & Acc 0,996, FPR 0,01753; sem métrica de eficácia da mitigação \\
Mazhar~\cite{mazhar2022ridps} & P, R, F1, taxa de deteção por n.º de \textit{bots} & P 100\%, R 0,97--0,98, F1 0,98--0,99; FP = 0 \\
Premakumari~\cite{premakumari2025qlearning} & Energia, PDR, latência, Acc de deteção, eficiência de mitigação & Deteção 93--98\%, mitigação 85--95\%; sem P, R, F1 nem FPR \\
Ramirez-Martinez~\cite{ramirezmartinez2025p4} & Acc, P, R, F1 & RF: Acc 97,02\% \textit{offline} e 98,28\% no \textit{testbed}; F1 97,21\% (Slowloris) \\
Rasheed~\cite{Rasheed2026} & P, R, F1, FPR & F1 0,900--0,956 e FPR 2,12--4,67\% nos quatro \textit{datasets} \\
Ravi~\cite{ravi2020semisupervised} & Acc, taxa de deteção, F, MCC & NSL-KDD: Acc 99,78\%, F 99,72\%, MCC 99,44\% \\
Rawat~\cite{rawat2024intelligent} & Acc, P, R, F1 & KNN: Acc 92,68\%, F1 0,85 \\
Rose~\cite{rose2022ideres} & Acc, P, R, F1, FPR & CNN: Acc 98,35\%, F1 99,16\%, FPR 0,98\%; 13 cenários mitigados \\
Saurabh~\cite{saurabh2022mlaps} & Acc, P, R, F1 & RF: Acc 98,4\%, F1 0,98 \\
Shukla~\cite{shukla2026cyber} & Acc, P, R, F1, FAR, macro-F1, AUROC, disponibilidade & Acc 97,46--98,64\%, F1 96,59--98,04\%, FAR 0,88--1,41\%; \textit{zero-day}: AUROC 0,963 \\
Srilekha~\cite{srilekha2025hybrid} & Acc, P, FPR, eficiência de resposta & Acc 99,2\%, FPR cerca de 1,5\% (lidos de gráficos) \\
Tariq~\cite{tariq2025saegru} & Acc, P, R, F1, AUC, validação cruzada (10 partes) & Acc 98,65\%, F1 98,27\%, AUC 0,993; FPR só em gráfico \\
Tawfik~\cite{tawfik2021internet} & Acc, F1 & LSTM: Acc 98,3\%, F1 85,9\% (binário) \\
Toony~\cite{toony2024multiblock} & Acc, P, R, F1, AUC, MCC, FPR, AME & Acc 99,75\% (Edge-IIoTset), 98,17\% (ToN-IoT), 98,21\% (X-IIoTID); AME 88,29--99,95\%; FPR 0,40--5,88\% \\
Verma~\cite{verma2024petrak} & Acc, TPR, FPR, sensibilidade, especificidade & Acc 99,70--99,998\% (três \textit{datasets}), 98,8\% no ToN-IoT; FPR 0,000--0,003 \\
Zolotukhin~\cite{zolotukhin2021curious} & TPR, FPR e recompensa (só em curvas) & Sem valores numéricos \\
\bottomrule
\end{tabular}

\vspace{3mm}
\caption{Tipo de avaliação, ambiente e custo computacional (QF4)}\label{tab:qf4}
\centering\scriptsize
\setlength{\tabcolsep}{3pt}
\renewcommand{\arraystretch}{0.84}
\begin{tabular}{@{}>{\raggedright\arraybackslash}p{0.13\textwidth}>{\raggedright\arraybackslash}p{0.14\textwidth}>{\raggedright\arraybackslash}p{0.30\textwidth}>{\raggedright\arraybackslash}p{0.39\textwidth}@{}}
\toprule
Estudo & Avaliação & Ambiente / \textit{hardware} & Custo e tempos reportados \\
\midrule
Kponyo~\cite{kponyo2020lightweight} & Tempo real & 4 Raspberry Pi 3B (dispositivos) & Deteção 0,10\,s; mitigação 0,22\,s; 21--28\,MB; +35,8\,mW \\
Ravi~\cite{ravi2020semisupervised} & \textit{Testbed} real & Cisco Nexus (\textit{fog}), 30 Raspberry Pi & Treino 1986\,s; teste 6,6\,s; resposta $\ge$34\% mais rápida que na \textit{cloud} \\
Dange~\cite{dange2025dual} & \textit{Testbed} real (ML \textit{offline}) & Raspberry Pi 4, ThingsBoard & Só tempo de treino \\
Karam~\cite{karam2025isaaf} & \textit{Hardware} real e \textit{offline} & Raspberry Pi 4 (\textit{edge}), \textit{broker} MQTT, ThingsBoard & Deteção $\approx$17\,s; mitigação 1--3\,s; CPU do \textit{broker} 60--80\% sob DoS \\
Srilekha~\cite{srilekha2025hybrid} & \textit{Hardware} real e simulação & Raspberry Pi 4, AWS EC2, Hyperledger Fabric & Resposta $<$3\,s; \textit{blockchain} no Raspberry Pi: +15\% de CPU, +50\,MB \\
Garba~\cite{garba2024sdn} & \textit{Hardware} real & Raspberry Pi (OpenFlow), Ryu & Só tempos de treino e teste (2\,min a 1\,h\,40) \\
Iqbal~\cite{Iqbal2026} & \textit{Hardware} real e \textit{offline} & Raspberry Pi 5 (10\,W); \textit{desktop} i7 & Latência $<$30\,ms \\
Al Abdulwahid~\cite{Abdulwahid2025} & Tempo real (servidor) & i9, 64\,GB, RTX 3090 & Latência 12,8\,ms; 275\,MB \\
Andrade~\cite{andrade2024ddos} & \textit{Testbed} 5G & Open5GS, UERANSIM, P4, ONOS & Deteção 20--140\,s; mitigação 2--4\,s; CPU 2--60\% \\
Rose~\cite{rose2022ideres} & \textit{Testbed} virtual & VMs (Cyber-Trust) & Só tempo de treino \\
Ramirez-Martinez~\cite{ramirezmartinez2025p4} & Emulação & Mininet, BMv2, VM 4 \textit{cores} & Deteção 20,6\,s; mitigação 0,128\,s \\
El-Sayed~\cite{elsayed2024mpguard} & Emulação & Mininet-WiFi, P4 & Deteção 11--33\,ms; CPU 20--35\% \\
Zolotukhin~\cite{zolotukhin2021curious} & Emulação & VMs, OVS, OpenDaylight & Não reportado \\
Tawfik~\cite{tawfik2021internet} & Emulação e \textit{offline} & Mininet e Ryu num Raspberry Pi 3B; IoTID20 & Não reportado \\
Toony~\cite{toony2024multiblock} & Emulação & Mininet-WiFi, BMv2, Ryu; VM i7 & Deteção 2--19\,ms; mitigação $\approx$3\,ms; CPU 100\% sob ataque \\
Gowtham~\cite{gowtham2026fedmp} & Emulação & Mininet SDN-\textit{fog}; i5 e 20 Raspberry Pi 4 & Inferência 196\,ms sob ataque; CPU do \textit{fog} 23\%; controlo 42--51\,ms \\
Mazhar~\cite{mazhar2022ridps} & Emulação & Mininet-WiFi, ONOS & Agente no dispositivo: +0,4\,pp de CPU, +0,3\,pp de memória \\
Tariq~\cite{tariq2025saegru} & Emulação e \textit{offline} & Servidores Xeon; nós \textit{edge} com GPU & Não reportado em valores \\
Ibrahim~\cite{ibrahim2026purechain} & Simulação & Portátil i5, 4\,GB & Inferência 1,6\,ms; CPU 14,5\%; 448\,MB \\
Shukla~\cite{shukla2026cyber} & Simulação & SUMO; RTX 3090 & Resposta 34,9\,ms (deteção 18,6; mitigação 16,3); 6,4\,GB de GPU \\
Verma~\cite{verma2024petrak} & Simulação e \textit{offline} & NS-3; Weka num i7 & Só tempos de construção e teste do modelo \\
Premakumari~\cite{premakumari2025qlearning} & Simulação & MATLAB Simulink, 50 nós & Latência 120\,ms; energia 45--74\,J \\
Karibandi~\cite{karibandi2025cyberdome} & Simulação e protótipo & \textit{Cloud}, \textit{gateways} e SDN & Inferência 8,74\,µs por amostra; treino 39,6\,s \\
Hunde~\cite{hunde2026rankshield} & Simulação & Cooja (Zolertia Z1), 19--100 nós & Treino 218--2750\,ms \\
Saurabh~\cite{saurabh2022mlaps} & \textit{Offline} & UNSW-NB15 & Tempo por modelo (DT 4\,ms; RF 46\,ms) \\
Rawat~\cite{rawat2024intelligent} & \textit{Offline} & IoT-23 & Não reportado \\
Hussain~\cite{hussain2026postquantum} & \textit{Offline} & Google Colab & \textit{Handshake} pós-quântico 3,1--4,4\,s \\
Kumar e Kumar~\cite{kumar2025advanced} & \textit{Offline} & PC i3, 8\,GB & Treino 250\,s; teste 15\,s \\
Rasheed~\cite{Rasheed2026} & \textit{Offline} & GPU A100 & Latência 15,6\,ms; 2,4\,GB de RAM \\
Arulkumar~\cite{Arulkumar2026} & \textit{Offline} & i7, 16\,GB, RTX & Latência 200\,ms \\
Chaudhary~\cite{Chaudhary2025} & \textit{Offline} & i5 e Azure & 115--170\,s por modelo \\
\bottomrule
\end{tabular}
\end{table*}
Esta secção responde à QF2 com base nos 31 estudos incluídos. Os valores
apresentados são os reportados por cada artigo.

\textbf{\textit{Datasets} utilizados.} A Tabela~\ref{tab:datasets} detalha os
\textit{datasets} de cada estudo. Muitos combinam \textit{datasets} públicos com dados
próprios, gerados em \textit{testbed}, emulação ou simulação: 17 dos 31 estudos usam
dados próprios, 13 usam apenas \textit{datasets} públicos e um não identifica o
\textit{dataset}~\cite{karibandi2025cyberdome}. O UNSW-NB15 e o IoT-23~\cite{iot23} são
os mais usados (cinco estudos cada), seguidos do ToN-IoT, do Edge-IIoTset e do CICIoT2023
(quatro cada). A dispersão é grande: os estudos usam 29 \textit{datasets} públicos
diferentes, 15 dos quais aparecem num único estudo.

Treze dos 31 estudos usam pelo menos um \textit{dataset} genérico de intrusão, que não é
de IoT, como o UNSW-NB15, o CICIDS2017 e 2018, o NSL-KDD, o KDDCUP'99, o CICDDoS2019 ou o
CTU-13, o que limita a sua representatividade para IoT. Os dados próprios raramente são
públicos: só o MQTTEEB-D~\cite{karam2025isaaf} e o \textit{dataset} do
IDERES~\cite{rose2022ideres} estão num repositório aberto, e o do SAE-GRU apenas em
parte~\cite{tariq2025saegru}. Os restantes estão disponíveis a pedido ou não estão
disponíveis, o que dificulta a replicação. Vários estudos também não descrevem o tamanho
dos dados nem a divisão entre treino e
teste~\cite{kumar2025advanced,gowtham2026fedmp,tariq2025saegru,toony2024multiblock}.


\textbf{Tipos de ataque detetados.} O DoS/DDoS está presente em 28 dos 31 estudos, e em
sete é o único
alvo~\cite{gowtham2026fedmp,mazhar2022ridps,garba2024sdn,ramirezmartinez2025p4,andrade2024ddos,kponyo2020lightweight,Chaudhary2025}.
Seguem-se as \textit{botnets} e a comunicação de C\&C (15 estudos), o \textit{malware} (13)
e a força bruta (9). A Tabela~\ref{tab:ataques} mostra a distribuição por estudo.

Os ataques detetados não coincidem com aqueles contra os quais a resposta automática é
avaliada. Vários estudos listam muitos tipos de ataque porque o \textit{dataset} os
contém, mas demonstram a resposta sobretudo contra \textit{floods} SYN, UDP, ICMP ou HTTP.
A deteção de ataques \textit{zero-day} só é testada explicitamente em dois estudos:
Shukla e Agarwal obtêm uma AUROC de 0,963~\cite{shukla2026cyber}, e no IDERES só três de
16 ataques \textit{zero-day} são mitigados~\cite{rose2022ideres}. Gowtham
\textit{et al.} limitam-se a invocá-la~\cite{gowtham2026fedmp}.


\textbf{Métricas de avaliação e resultados.} As métricas usadas derivam da matriz de
confusão e estão definidas nas Eqs.~(\ref{eq:precrec})--(\ref{eq:fpr}). A
\textit{accuracy} é reportada em 28 dos 31 estudos e o F1 em 25, mas só 13 dão um valor
explícito da taxa de falsos
positivos~\cite{srilekha2025hybrid,shukla2026cyber,kumar2025advanced,gowtham2026fedmp,toony2024multiblock,verma2024petrak,rose2022ideres,elsayed2024mpguard,Abdulwahid2025,Iqbal2026,Rasheed2026,Arulkumar2026,Chaudhary2025},
e poucos medem a eficácia da resposta. Os valores da Tabela~\ref{tab:metricas} são os que
cada estudo reporta e não são comparáveis entre si, porque os \textit{datasets}, as
divisões e as tarefas (binária ou multiclasse) diferem.


As métricas de classificação concentram-se acima de 97\%, mas a robustez é pouco
demonstrada. Quando Karam \textit{et al.} testam o modelo noutro \textit{dataset}, a
\textit{accuracy} cai para 8--21\%~\cite{karam2025isaaf}; Hussain \textit{et al.} obtêm um
F1 de 0,751 no NSL-KDD, contra cerca de 1,0 nos restantes
\textit{datasets}~\cite{hussain2026postquantum}; e no trabalho de Dange e Nitnaware a
\textit{accuracy} passa de 99,97\% no BoT-IoT para 92\% no
\textit{testbed}~\cite{dange2025dual}.

A eficácia da resposta, central para a questão geral, só tem métrica própria em nove
dos 31 estudos: a AME de Toony \textit{et al.}~\cite{toony2024multiblock}, a eficiência de
resposta de Srilekha \textit{et al.}~\cite{srilekha2025hybrid}, a taxa de mitigação de
Gowtham \textit{et al.}, de 91,9 a 94,5\%~\cite{gowtham2026fedmp}, e de Premakumari
\textit{et al.}~\cite{premakumari2025qlearning}, a disponibilidade sob ataque de Shukla e
Agarwal~\cite{shukla2026cyber}, a taxa de sucesso dos ataques de Ibrahim
\textit{et al.}~\cite{ibrahim2026purechain}, a recompensa de Zolotukhin \textit{et al.},
que soma o tráfego legítimo que passa e o malicioso que é
bloqueado~\cite{zolotukhin2021curious}, os cenários mitigados do
IDERES~\cite{rose2022ideres} e o efeito na entrega de pacotes e na latência do
RankShield~\cite{hunde2026rankshield}. Kumar e Kumar, por exemplo, propõem uma mitigação mas não
medem a sua eficácia~\cite{kumar2025advanced}. Foram ainda identificadas incoerências
internas em alguns estudos, entre o resumo e as
tabelas~\cite{hussain2026postquantum,premakumari2025qlearning,kumar2025advanced} ou em
contagens que mudam ao longo do texto~\cite{mazhar2022ridps}, o que constitui uma
limitação das fontes.

\textbf{Síntese e lacunas.} Os resultados reportados são muito altos, mas assentam em
\textit{datasets} e métricas pouco comparáveis e raramente medem a contenção. Destacam-se
quatro lacunas. A primeira é que os \textit{datasets} são pouco representativos: cerca de
40\% dos estudos usam \textit{datasets} genéricos de intrusão, e os dados próprios são
quase sempre simulados ou emulados e não são públicos. A segunda é que a generalização
raramente é testada fora do \textit{dataset} de treino, e quando o é os resultados caem.
A terceira é que os ataques se concentram em \textit{floods}: o DoS/DDoS domina, a
resposta avaliada é quase sempre contra \textit{floods}, e as \textit{botnets} com C\&C e
a força bruta têm pouca cobertura. A quarta é que as métricas de resposta são escassas:
menos de um terço dos estudos mede a eficácia da contenção ou o seu efeito no tráfego
legítimo, e menos de metade dá um valor explícito da taxa de falsos positivos. Por isso a
solução proposta reporta a FPR e mede as quarentenas indevidas com tráfego só benigno
(cenário S1, Secção~\ref{sec:proposta}).

\subsection{Mecanismos de resposta automática (QF3)}
% RESPONSÁVEL: Frente 5 (Gonçalo Batista)
A análise dos estudos selecionados permitiu identificar diferentes abordagens para a resposta automática a ataques em redes IoT. Os mecanismos encontrados incluem o descarte de pacotes, o bloqueio de fluxos, o isolamento de dispositivos e a reconfiguração dinâmica de políticas de segurança. As principais diferenças entre os estudos estão no grau de intervenção, no componente responsável pela execução da resposta e nas condições que levam à sua ativação.

O bloqueio de tráfego é uma das medidas mais diretas. Kponyo \textit{et al.}~\cite{kponyo2020lightweight} apresentam uma solução que funciona no próprio dispositivo IoT, utilizando um classificador k-NN para identificar ataques DoS e o NetfilterQueue para descartar os pacotes associados à ameaça. Neste caso, a deteção e a resposta são realizadas localmente, sem depender da intervenção de um controlador central. Em contraste, o MP-GUARD, de El-Sayed \textit{et al.}~\cite{elsayed2024mpguard}, combina aprendizagem automática com uma arquitetura SDN de múltiplos controladores e \textit{switches} P4. A identificação de tráfego suspeito permite desencadear políticas de filtragem, limitação de débito e isolamento de recursos na infraestrutura de rede. Estas abordagens mostram que o bloqueio pode ser aplicado tanto no equipamento protegido como em pontos intermédios por onde passam as suas comunicações.

Uma estratégia mais abrangente é apresentada pelo CyberDome, de Karibandi \textit{et al.}~\cite{karibandi2025cyberdome}. O sistema combina um IDS baseado em XGBoost, executado na \textit{cloud}, com um módulo de mitigação baseado em SDN. A resposta é organizada em três níveis: controlo dinâmico de fluxos, segmentação da rede e quarentena do dispositivo. Quando é detetada atividade maliciosa, o sistema pode começar por restringir o tráfego e aumentar progressivamente o nível de isolamento caso a ameaça persista. Esta abordagem distingue-se por permitir aplicar medidas com diferentes níveis de severidade, em vez de recorrer imediatamente ao bloqueio total. Embora os autores demonstrem os diferentes níveis de mitigação numa rede simulada, a avaliação quantitativa concentra-se sobretudo no desempenho do classificador, não permitindo determinar com o mesmo detalhe a eficácia e a latência das medidas de contenção.

O RankShield, proposto por Hunde e Ismail~\cite{hunde2026rankshield}, apresenta outro mecanismo de isolamento, direcionado para ataques ao protocolo de encaminhamento RPL. O sistema utiliza uma rede neuronal profunda (DNN) para identificar alterações anormais dos valores de rank anunciados pelos dispositivos. Antes de executar a resposta, essa previsão é sujeita a uma verificação adicional baseada em limiares dinâmicos, procurando distinguir ataques de alterações legítimas na topologia. Quando a atividade maliciosa é confirmada, o nó atacante é colocado numa lista negra, cuja informação é propagada pela rede para impedir a sua participação no encaminhamento. Ao contrário de soluções que dependem de \textit{firewalls} ou controladores SDN, o RankShield aplica a contenção através dos próprios mecanismos de comunicação da rede RPL. A validação adicional constitui também uma forma de reduzir o risco de isolamento indevido.

Estas diferenças evidenciam que o local de aplicação da resposta depende da arquitetura e do tipo de ameaça. A solução de Kponyo \textit{et al.}~\cite{kponyo2020lightweight} atua no próprio dispositivo, enquanto o MP-GUARD~\cite{elsayed2024mpguard} utiliza controladores e \textit{switches} programáveis. O CyberDome~\cite{karibandi2025cyberdome} centraliza a deteção e a gestão na \textit{cloud}, mas recorre à infraestrutura SDN para aplicar as restrições. Já o RankShield~\cite{hunde2026rankshield} adota uma abordagem distribuída, em que os nós da rede participam no isolamento dos atacantes. As soluções centralizadas facilitam a gestão de políticas para vários equipamentos, ao passo que as abordagens locais ou distribuídas podem reduzir a dependência de um único componente. Contudo, estas últimas exigem que os dispositivos envolvidos tenham capacidade para executar os mecanismos de proteção.

Relativamente ao desencadeamento da resposta, os estudos apresentam estratégias que vão desde a aplicação de regras predefinidas até à seleção dinâmica de contramedidas. No IDERES, Rose \textit{et al.}~\cite{rose2022ideres} combinam deteção de intrusões com grafos de ataques e um modelo de decisão probabilístico baseado em POMDP (Partially Observable Markov Decision Process). Este mecanismo considera os alertas recebidos, a informação disponível sobre a rede e os possíveis caminhos de ataque para selecionar regras de \textit{firewall}. A resposta pode consistir no bloqueio de uma comunicação específica ou numa restrição mais abrangente das ligações de um equipamento. Desta forma, procura-se evitar que todas as ameaças desencadeiem automaticamente o mesmo tipo de bloqueio, independentemente do seu contexto.

Uma abordagem adaptativa é também apresentada no Curious SDN, de Zolotukhin \textit{et al.}~\cite{zolotukhin2021curious}, que utiliza aprendizagem por reforço para modificar dinamicamente as políticas de segurança. O agente observa o estado da rede e pode tomar decisões como alterar regras de fluxo, redirecionar tráfego ou ajustar a sensibilidade dos mecanismos de deteção. A aprendizagem procura equilibrar a quantidade de tráfego malicioso bloqueado com a preservação das comunicações legítimas. No entanto, os resultados experimentais mostram que o desempenho da política aprendida depende do cenário, por exemplo, a abordagem apresenta resultados favoráveis em ataques DDoS e \textit{botnets}, mas revela dificuldades em determinados ataques a aplicações \textit{web}. Isto demonstra que a utilização de mecanismos inteligentes não garante, por si só, respostas eficazes em todas as situações.

A validação experimental dos mecanismos de contenção constitui uma limitação importante na comparação dos estudos. O MP-GUARD~\cite{elsayed2024mpguard} e o Curious SDN~\cite{zolotukhin2021curious} recorrem a ambientes de rede emulados, enquanto o RankShield~\cite{hunde2026rankshield} é avaliado através de simulações com o simulador Cooja. O IDERES~\cite{rose2022ideres} utiliza um testbed virtualizado e apresenta resultados de mitigação para diferentes cenários, mas a eficácia varia conforme o ataque. O trabalho de Kponyo \textit{et al.}~\cite{kponyo2020lightweight} distingue-se pela utilização de dispositivos Raspberry Pi, nos quais são efetivamente executadas ações de descarte de pacotes. Já o CyberDome~\cite{karibandi2025cyberdome} apresenta uma contribuição importante para a arquitetura da resposta, embora a avaliação da execução operacional da contenção seja menos aprofundada. A Secção~\ref{sec:qf4} aprofunda estas diferenças quanto aos ambientes de teste, tempos de resposta e custos computacionais.

Em conjunto, os estudos mostram que a deteção automática não deve ser confundida com a contenção efetiva. Um modelo pode apresentar uma elevada exatidão sem que isso demonstre que o ataque foi interrompido. Para avaliar a resposta, é necessário verificar se a medida foi executada, quanto tempo demorou a produzir efeito e se provocou restrições indevidas ao tráfego legítimo. Esta distinção é especialmente importante em redes IoT, onde o isolamento incorreto de um dispositivo pode comprometer serviços que dependem da sua disponibilidade.

% Visão de conjunto dos 31 estudos (acrescentado na integração, a partir da Tabela de comparação).
Alargando a análise aos 31 estudos incluídos, os mesmos padrões repetem-se
(Tabelas~\ref{tab:comparacao} e~\ref{tab:comparacao2}). Quanto ao local, em 12 estudos a
resposta é aplicada por controladores SDN ou \textit{switches}
P4~\cite{ramirezmartinez2025p4,andrade2024ddos,elsayed2024mpguard,zolotukhin2021curious,tawfik2021internet,Iqbal2026,Chaudhary2025,toony2024multiblock,gowtham2026fedmp,mazhar2022ridps,garba2024sdn,karibandi2025cyberdome},
em 12 num \textit{gateway} ou num nó \textit{edge} ou
\textit{fog}~\cite{rose2022ideres,ravi2020semisupervised,ibrahim2026purechain,dange2025dual,karam2025isaaf,srilekha2025hybrid,hussain2026postquantum,kumar2025advanced,tariq2025saegru,Abdulwahid2025,Arulkumar2026,verma2024petrak},
e em três nos próprios
dispositivos~\cite{kponyo2020lightweight,premakumari2025qlearning,hunde2026rankshield}; nos
restantes quatro o local não é especificado. Quanto à medida, o descarte ou o bloqueio de tráfego é o mais comum, e 11
estudos chegam a isolar ou a pôr em quarentena o dispositivo
inteiro~\cite{rose2022ideres,ibrahim2026purechain,dange2025dual,Iqbal2026,toony2024multiblock,srilekha2025hybrid,hussain2026postquantum,tariq2025saegru,karibandi2025cyberdome,hunde2026rankshield,shukla2026cyber}.
A limitação de débito surge em seis como medida
intermédia~\cite{elsayed2024mpguard,ibrahim2026purechain,srilekha2025hybrid,gowtham2026fedmp,kumar2025advanced,karibandi2025cyberdome},
e só dois tornam o bloqueio temporário, com um temporizador
aleatório~\cite{ravi2020semisupervised} ou uma quarentena de
120\,s~\cite{ibrahim2026purechain}. Quanto ao desencadeamento, vários estudos exigem
evidência acumulada ou uma segunda verificação antes de agir: um contador por IP que só
bloqueia depois de três fluxos maliciosos~\cite{ramirezmartinez2025p4}, cinco previsões
de ataque seguidas~\cite{andrade2024ddos}, limiares de
confiança~\cite{ibrahim2026purechain,gowtham2026fedmp,hussain2026postquantum}, a
confirmação do alerta por pelo menos dois nós~\cite{srilekha2025hybrid} ou limiares
dinâmicos~\cite{hunde2026rankshield}. Seis recorrem a modelos de decisão, sobretudo
aprendizagem por
reforço~\cite{rose2022ideres,zolotukhin2021curious,gowtham2026fedmp,premakumari2025qlearning,kumar2025advanced,Arulkumar2026}.
Em quatro estudos, pelo contrário, a resposta é proposta mas não chega a ser
implementada, ou fica-se pelo plano
conceptual~\cite{rawat2024intelligent,hussain2026postquantum,kumar2025advanced,shukla2026cyber}.

Em suma, e respondendo à QF3, os mecanismos de resposta automática identificados na literatura consistem principalmente no descarte e bloqueio de tráfego, na aplicação de regras de \textit{firewall}, no isolamento de nós ou dispositivos e na reconfiguração de políticas através de SDN. As ações são aplicadas em diferentes níveis, desde o próprio dispositivo IoT até aos \textit{gateways} e equipamentos programáveis, podendo a decisão ser realizada localmente ou em componentes centralizados, incluindo a \textit{cloud}. O desencadeamento depende de classificações produzidas por modelos de ML, de mecanismos adicionais de validação ou de modelos de decisão adaptativa. Apesar da diversidade de soluções, a revisão evidencia que a eficácia da contenção continua a ser avaliada de forma menos consistente do que a deteção, reforçando a importância de desenvolver e testar sistemas que integrem ambas as etapas.

\subsection{Avaliação das soluções (QF4)}\label{sec:qf4}
% RESPONSÁVEL: Frente 4 (Rafael Reis)
A literatura mostra uma evolução clara: muitos trabalhos começam com validações em
ambiente controlado, \textit{offline} e com \textit{datasets} de referência, e evoluem
para implementações em protótipos, \textit{testbeds} e redes em tempo real. As abordagens
de avaliação dividem-se, assim, em duas categorias.

\textbf{Avaliação estática ou \textit{offline}.} É a etapa em que se calculam as
métricas (exatidão, taxa de falsos positivos, matrizes de confusão) e se ajustam os
parâmetros dos modelos, e em sete dos 31 estudos é a única. O MLAPS é avaliado apenas sobre o
UNSW-NB15~\cite{saurabh2022mlaps}; Rawat \textit{et al.} classificam fluxos do IoT-23 e,
embora o resumo refira simulações e implementações reais, os resultados apresentados são
apenas de classificação~\cite{rawat2024intelligent}. O mesmo acontece em mais
cinco estudos~\cite{hussain2026postquantum,kumar2025advanced,Rasheed2026,Arulkumar2026,Chaudhary2025}.

\textbf{Avaliação em tempo real, protótipos e \textit{testbeds}.} A literatura orientada
à aplicação prática valida os modelos em cenários simulados ou reais, com graus de
realismo muito diferentes (Tabela~\ref{tab:qf4}). Seis estudos usam
simulação~\cite{ibrahim2026purechain,shukla2026cyber,verma2024petrak,premakumari2025qlearning,hunde2026rankshield,karibandi2025cyberdome}
e oito emulação de rede, quase sempre com
Mininet~\cite{ramirezmartinez2025p4,elsayed2024mpguard,zolotukhin2021curious,tawfik2021internet,toony2024multiblock,gowtham2026fedmp,mazhar2022ridps,tariq2025saegru}.
Tawfik \textit{et al.}, por exemplo, correm o Mininet e o controlador Ryu num Raspberry
Pi 3B, mas deixam a aprendizagem profunda para uma avaliação \textit{offline} sobre o
IoTID20~\cite{tawfik2021internet}. Dois usam um \textit{testbed} virtualizado ou uma rede
5G em \textit{software}~\cite{rose2022ideres,andrade2024ddos}, e um corre em tempo real
num servidor~\cite{Abdulwahid2025}. Só sete correm em \textit{hardware} físico, quase
sempre com Raspberry
Pi~\cite{ravi2020semisupervised,dange2025dual,kponyo2020lightweight,karam2025isaaf,srilekha2025hybrid,garba2024sdn,Iqbal2026}.
O caso mais próximo de um dispositivo real é o de Kponyo \textit{et al.}, em que quatro
Raspberry Pi 3B fazem de dispositivos IoT e executam eles próprios a deteção e a
resposta~\cite{kponyo2020lightweight}. Mesmo nestes estudos, o treino e a avaliação dos
modelos são normalmente feitos \textit{offline}, e o protótipo serve para demonstrar a
resposta.


O tempo até à contenção é medido em poucos estudos, e os resultados mostram onde está o
estrangulamento. No caso do Slowloris, a mitigação demora 0,128\,s, mas a deteção leva em
média 20,6\,s, porque um ataque lento só se revela ao fim de algum
tempo~\cite{ramirezmartinez2025p4}. Na rede 5G, a mitigação demora 2 a 4\,s e a deteção
entre 20 e 140\,s, conforme o intervalo de recolha das
estatísticas~\cite{andrade2024ddos}, e no ISAAF a deteção leva cerca de 17\,s, contra 1
a 3\,s do bloqueio~\cite{karam2025isaaf}. Nestes casos, o tempo de deteção é uma a duas
ordens de grandeza maior do que o da resposta, o que justifica medir separadamente as duas
parcelas de $T_c$ (Eq.~\ref{eq:tc}). Os tempos de milissegundos aparecem apenas em
emulação ou simulação~\cite{toony2024multiblock,shukla2026cyber}, e o estudo de Kponyo
\textit{et al.} é a exceção em \textit{hardware} real, com deteção em 0,10\,s e mitigação em
0,22\,s~\cite{kponyo2020lightweight}, mas trata \textit{floods}, que se revelam muito
depressa.

O custo computacional é reportado de forma desigual. Os valores mais completos vêm de
estudos que correm em equipamentos de baixo custo. Kponyo \textit{et al.} medem a memória
ocupada pelo mecanismo (21 a 28\,MB, com um acréscimo de 2,82 pontos percentuais na
ocupação de memória de um dos nós) e um consumo adicional de 35,8\,mW, embora descrevam a
CPU apenas de forma qualitativa~\cite{kponyo2020lightweight}; no trabalho de Srilekha
\textit{et al.}, o cliente de \textit{blockchain} acrescenta cerca de 15\% de CPU e 50\,MB
de memória a um Raspberry Pi~\cite{srilekha2025hybrid}; e no R-IDPS o agente instalado nos
dispositivos acrescenta menos de meio ponto percentual de CPU e de
memória~\cite{mazhar2022ridps}. Os restantes reportam tempos de inferência, a utilização
de CPU dos controladores, entre 2\% e 35\%~\cite{andrade2024ddos,elsayed2024mpguard}, ou
apenas tempos de treino, que no caso de modelos complexos treinados sobre milhões de
registos podem ir de alguns minutos a quase duas horas~\cite{garba2024sdn}. Quatro estudos
não reportam qualquer
custo~\cite{zolotukhin2021curious,tawfik2021internet,rawat2024intelligent,tariq2025saegru}.

Em síntese, a passagem da avaliação \textit{offline} para protótipos em \textit{hardware}
real mostra que a deteção e a resposta podem correr em equipamentos de baixo custo com
custos reduzidos~\cite{kponyo2020lightweight,srilekha2025hybrid,Iqbal2026}. Essa passagem,
porém, só foi feita em sete dos 31 estudos, e quase sempre com um único tipo de ataque.
É esta lacuna que a avaliação da solução proposta procura cobrir, ao medir $T_c$, a CPU e
a memória do \textit{gateway} nos mesmos cenários para os três modelos.


\begin{table*}[p]
\caption{Comparação dos trabalhos relacionados (S: supervisionado; NS: não supervisionado; SS: semi-supervisionado; R: por reforço)}\label{tab:comparacao}
\centering\scriptsize
\setlength{\tabcolsep}{2.5pt}
\renewcommand{\arraystretch}{0.97}
\begin{tabular}{@{}>{\raggedright\arraybackslash}p{0.10\textwidth}>{\raggedright\arraybackslash}p{0.12\textwidth}>{\raggedright\arraybackslash}p{0.14\textwidth}c>{\raggedright\arraybackslash}p{0.11\textwidth}>{\raggedright\arraybackslash}p{0.15\textwidth}>{\raggedright\arraybackslash}p{0.10\textwidth}>{\raggedright\arraybackslash}p{0.115\textwidth}c@{}}
\toprule
Referência & \textit{Dataset}(s) & Modelo(s) & Apr. & Onde atua a resposta & Resposta automática & Avaliação & Contenção medida & QA \\
\midrule
Noetzold \textit{et al.}~\cite{noetzold2026anomaly} (2026) & CICIDS2017, CVE, DNS e texto (Kaggle) & BiLSTM, FFNN, CNN, RNN-CNN & S & -- & Só alertas, numa interface em tempo real & \textit{Offline}; inferência em servidor, Raspberry Pi 4 e ESP32-S3 & \nao & -- \\
\midrule
Kponyo \textit{et al.}~\cite{kponyo2020lightweight} (2020) & Próprio (CPU e memória) & KNN sobre CPU e memória & S & Dispositivo (\textit{host}) & Descarte por protocolo e porta (NetfilterQueue) & Tempo real (4 Raspberry Pi 3B) & Mitigação em 0,22\,s & 7,0 \\
Ravi e Shalinie~\cite{ravi2020semisupervised} (2020) & NSL-KDD, próprio & DFNN + K-means (SDRK) & SS & Nó \textit{fog} (\textit{switch}) & ACL de descarte do par MAC--IP com temporizador & \textit{Testbed} real (30 Raspberry Pi) & Resposta $\ge$34\% mais rápida que na \textit{cloud} & 6,5 \\
Tawfik \textit{et al.}~\cite{tawfik2021internet} (2021) & IoTID20, próprio & RF, CNN, LSTM & S & Controlador SDN (Ryu) & Descarte ou bloqueio da porta ou do \textit{host} (descrito) & Emulação (Mininet) e \textit{offline} & \nao & 6,0 \\
Zolotukhin \textit{et al.}~\cite{zolotukhin2021curious} (2021) & CICIDS2018 & PPO + MLP, CNN, LSTM, AE, SOM, BiGAN & S, NS, R & Controlador SDN & Descarte e redirecionamento de fluxos & Emulação & Curvas de recompensa & 5,0 \\
Mazhar \textit{et al.}~\cite{mazhar2022ridps} (2022) & Próprio & SVM + assinaturas & S & Controlador SDN (ONOS) & Regras OpenFlow de descarte & Emulação (Mininet-WiFi) & \nao{} (atraso do tráfego legítimo: 8\,s) & 6,5 \\
Rose \textit{et al.}~\cite{rose2022ideres} (2022) & Próprio (público) & CNN + grafos de ataque (POMDP) & S & \textit{Gateway} e \textit{cloud} & Regras iptables: ligação específica ou isolamento do \textit{host} & \textit{Testbed} virtualizado & Cenários mitigados (sem tempo) & 6,5 \\
Saurabh \textit{et al.}~\cite{saurabh2022mlaps} (2022) & UNSW-NB15 & DT + RF (2 níveis) & S & Não especificado & Descarte de pacotes, aplicado ao \textit{dataset} & \textit{Offline} & \nao & 6,0 \\
Andrade \textit{et al.}~\cite{andrade2024ddos} (2024) & CICDDoS2019 & LR + NB & S & \textit{Switch} P4 (ONOS) & IP do equipamento em tabela de descarte após 5 previsões & \textit{Testbed} 5G & Mitigação em 2--4\,s & 6,5 \\
El-Sayed \textit{et al.}~\cite{elsayed2024mpguard} (2024) & ToN-IoT, Edge-IIoTset, CICIoT2023 & \textit{Stacking} (RF, SVM, DT, kNN, \dots{} + MLP) & S & \textit{Switches} P4 e controladores SDN & Descarte e limitação de débito & Emulação (Mininet-WiFi) & \nao{} (só deteção: 11--33\,ms) & 5,5 \\
Garba \textit{et al.}~\cite{garba2024sdn} (2024) & UNSW-NB15, CICDDoS2019, próprio & DT, KNN, LR, SVM & S & Controlador SDN (Ryu) & Descarte do MAC atacante (desencadeado pelo Snort) & \textit{Hardware} real (Raspberry Pi) & \nao & 6,0 \\
Rawat \textit{et al.}~\cite{rawat2024intelligent} (2024) & IoT-23 & KNN, LR, LSTM & S & Não especificado & Contramedida proposta, não implementada & \textit{Offline} & \nao & 5,0 \\
Toony \textit{et al.}~\cite{toony2024multiblock} (2024) & Edge-IIoTset, ToN-IoT, X-IIoTID & \textit{Ensemble} ponderado de 6 classificadores (EWEA) & S & \textit{Switches} P4 e controladores SDN & Descarte da origem, isolamento do nó e alerta aos outros controladores & Emulação (Mininet-WiFi) & Mitigação em $\approx$3\,ms & 7,5 \\
Verma \textit{et al.}~\cite{verma2024petrak} (2024) & KDDCUP'99, ToN-IoT, \dots, próprio & BayesNet, regressão logística & S & RSU (\textit{edge} veicular) & Desvio (\textit{sinkhole}) e descarte por assinatura & Simulação (NS-3) e \textit{offline} & \nao & 6,0 \\
Al Abdulwahid~\cite{Abdulwahid2025} (2025) & CICIDS2017, CSE-CIC-IDS2018, UNSW-NB15 & AE + RNN (CyberGuard-X) & SS & \textit{Gateway} ou \textit{edge} & Automática (tipo não indicado na extração) & Tempo real (i9, RTX 3090) & Latência 12,8\,ms & 7,5 \\
Chaudhary \textit{et al.}~\cite{Chaudhary2025} (2025) & BoT-IoT & SVM, RF, J48, KNN, XGBoost & S & \textit{Fog} com SDN & Bloqueio & \textit{Offline} & \nao & 6,0 \\
\bottomrule
\end{tabular}
\end{table*}

\begin{table*}[p]
\caption{Comparação dos trabalhos relacionados (continuação)}\label{tab:comparacao2}
\centering\scriptsize
\setlength{\tabcolsep}{2.5pt}
\renewcommand{\arraystretch}{0.97}
\begin{tabular}{@{}>{\raggedright\arraybackslash}p{0.10\textwidth}>{\raggedright\arraybackslash}p{0.12\textwidth}>{\raggedright\arraybackslash}p{0.14\textwidth}c>{\raggedright\arraybackslash}p{0.11\textwidth}>{\raggedright\arraybackslash}p{0.15\textwidth}>{\raggedright\arraybackslash}p{0.10\textwidth}>{\raggedright\arraybackslash}p{0.115\textwidth}c@{}}
\toprule
Referência & \textit{Dataset}(s) & Modelo(s) & Apr. & Onde atua a resposta & Resposta automática & Avaliação & Contenção medida & QA \\
\midrule
Dange e Nitnaware~\cite{dange2025dual} (2025) & BoT-IoT, CTU-13, próprio & LR, SVM, NB + regras & S & \textit{Gateway} (Raspberry Pi) & Isolamento do dispositivo (não descrito) & \textit{Testbed} real & \nao & 5,0 \\
Karam \textit{et al.}~\cite{karam2025isaaf} (2025) & MQTTset, próprio (público) & DT, GRU (escolhidos entre 15 modelos) & S & \textit{Gateway} (Raspberry Pi 4) & Bloqueio do IP e da porta 1883 na \textit{firewall} do \textit{broker}, via REST & \textit{Hardware} real e \textit{offline} & Mitigação em 1--3\,s (deteção $\approx$17\,s) & 7,5 \\
Karibandi \textit{et al.}~\cite{karibandi2025cyberdome} (2025) & Não identificado (11 classes) & XGBoost & S & \textit{Cloud} e controlador SDN & Controlo de fluxos, segmentação e quarentena, por níveis & Simulação e protótipo & \nao & 5,0 \\
Kumar e Kumar~\cite{kumar2025advanced} (2025) & CICIoT2023 & DCSAtten-EBiGRU + Double DQN & S, R & \textit{Gateway} & Bloquear, permitir ou limitar (sem resultados) & \textit{Offline} & \nao & 5,0 \\
Premakumari \textit{et al.}~\cite{premakumari2025qlearning} (2025) & AWID, IoT-23, simulado & FNN + Q-learning e Double Q-learning & S, R & Dispositivo (nós sensores) & Ajuste do nível de encriptação (sem bloqueio) & Simulação (MATLAB) & Latência 120\,ms & 6,0 \\
Ramirez-Martinez \textit{et al.}~\cite{ramirezmartinez2025p4} (2025) & CICIoT2023, IoT-23, próprio & KNN, DT, RF, LSTM, CNN, GRU, MLP & S & \textit{Switch} P4 & Descarte do IP de origem após 3 fluxos maliciosos & Emulação (Mininet, BMv2) & Mitigação em 0,128\,s & 7,0 \\
Srilekha \textit{et al.}~\cite{srilekha2025hybrid} (2025) & UNSW-NB15, sintético & Extra Trees (\textit{edge}), CNN-LSTM (\textit{cloud}) & S & \textit{Edge} e \textit{cloud}, com \textit{blockchain} & Bloqueio, limitação de débito e quarentena por \textit{smart contract} & \textit{Hardware} real (Raspberry Pi 4) e simulação & Resposta $<$3\,s & 7,0 \\
Tariq e Ahanger~\cite{tariq2025saegru} (2025) & IoT-23, MedBIoT, próprio & SAE + GRU & S, NS & \textit{Gateway} ou \textit{edge} & Isolamento, bloqueio de IP e regras de \textit{firewall} por nível de risco & Emulação e \textit{offline} & \nao & 5,5 \\
Arulkumar e K~\cite{Arulkumar2026} (2026) & DAPT2020, KDDCUP'99, CICIoT2023, \dots & CNN-LSTM-BNN + XGBoost/LR; MDP multiagente & S, R & \textit{Gateway} ou \textit{edge} & Bloqueio e outras ações & \textit{Offline} & \nao{} (latência 200\,ms) & 7,0 \\
Gowtham \textit{et al.}~\cite{gowtham2026fedmp} (2026) & Edge-IIoTset, CIC-IDS2017, próprio & KPCA-TCN, Ghost-SE federado; Q-learning & S, R & Nós \textit{fog} e controlador SDN & Etiquetar, limitar, descartar ou reencaminhar fluxos & \textit{Testbed} Mininet SDN-\textit{fog} & Latência de controlo 42--51\,ms & 6,5 \\
Hunde e Ismail~\cite{hunde2026rankshield} (2026) & Próprio (Cooja) & DNN + limiares dinâmicos & S & Dispositivos (nós RPL) & \textit{Blacklist} do nó atacante, propagada pela rede & Simulação (Cooja) & \nao{} (PDR +67\%, latência $-$66\%) & 5,5 \\
Hussain \textit{et al.}~\cite{hussain2026postquantum} (2026) & NSL-KDD, MedBIoT, CIC-IDS2017, \dots & RF, LR, MLP (\textit{trust score}) & S & \textit{Gateway} e motor de política & \textit{Allow}, \textit{monitor}, \textit{deny} ou quarentena (não implementadas) & \textit{Offline} (Colab) & \nao & 7,0 \\
Ibrahim \textit{et al.}~\cite{ibrahim2026purechain} (2026) & IoTForge Pro, WUSTL-IIoT, X-IIoTID & CNN-BiLSTM & S & \textit{Gateway} (não especificado) & Quarentena de 120\,s ou limitação de débito & Simulação (\textit{replay}) & Taxa de sucesso dos ataques 0--8,8\% & 6,0 \\
Iqbal \textit{et al.}~\cite{Iqbal2026} (2026) & WUSTL-IIoT-2021, X-IIoTID & RF + MLP (DeepGuard) & S & Controlador SDN & Bloqueio, quarentena em VLAN e contramedidas por domínio & \textit{Hardware} real (Raspberry Pi 5) e \textit{offline} & Latência $<$30\,ms & 7,0 \\
Rasheed \textit{et al.}~\cite{Rasheed2026} (2026) & SWaT, WADI, BATADAL, IIoT-Security & TAC-LSTM-GNN + ARLM & S & -- & Bloqueio adaptado ao tipo de ataque & \textit{Offline} & \nao{} (latência 15,6\,ms) & 6,0 \\
Shukla e Agarwal~\cite{shukla2026cyber} (2026) & UNSW-NB15, CICIDS2018, BoT-IoT & ST-GT (\textit{graph transformer}) & S, SS & Não especificado & Isolamento e controlo de acesso (só conceptual) & Simulação (SUMO) & 34,9\,ms (simulação) & 7,0 \\
\midrule
\textbf{Este trabalho} & IoT-23 & RF, IF, AE & S, NS, SS & \textit{Gateway} (nftables) & Quarentena do dispositivo após $k$ de $w$ fluxos, com expiração & Protótipo em contentores & $T_c$, CPU e memória (planeado) & -- \\
\bottomrule
\end{tabular}
\end{table*}
\clearpage
\subsection{Comparação com o estado da arte}
% RESPONSÁVEL: Frente 1 (Afonso) — vale 9 dos 18 valores do relatório!

% Uma linha por estudo incluído, por ordem cronológica.
A Tabela~\ref{tab:comparacao} compara os estudos incluídos segundo os critérios do
formulário de extração, mais o trabalho de Noetzold \textit{et al.}~\cite{noetzold2026anomaly},
que serviu de ponto de partida, e a solução proposta neste trabalho. As colunas
correspondem às perguntas de pesquisa: modelos e tipo de aprendizagem (QF1),
\textit{datasets} (QF2), mecanismo de resposta e local onde atua (QF3), tipo de avaliação
e medição da contenção (QF4).

\textbf{Comparação com revisões anteriores.} Antes de comparar os estudos, importa
situar esta revisão face às que já existem. A Tabela~\ref{tab:revisoes} compara-a com as
quatro revisões usadas para afinar as perguntas de pesquisa
(Secção~\ref{sec:metodologia}). Todas tratam a deteção e algumas referem a mitigação, mas
nenhuma tem a resposta automática como objeto central, nem analisa onde é aplicada, como
é desencadeada ou se chega a ser avaliada. Esta revisão faz das três questões as suas
perguntas QF3 e QF4 e usa o EC6 para reter apenas estudos que abordam a resposta.

\begin{table}[H]
\caption{Comparação desta revisão com revisões anteriores (n.i.: não indicado no resumo)}\label{tab:revisoes}
\centering\scriptsize
\setlength{\tabcolsep}{2.5pt}
\begin{tabular}{@{}>{\raggedright\arraybackslash}p{0.16\columnwidth}>{\raggedright\arraybackslash}p{0.24\columnwidth}>{\raggedright\arraybackslash}p{0.12\columnwidth}>{\raggedright\arraybackslash}p{0.13\columnwidth}>{\raggedright\arraybackslash}p{0.11\columnwidth}>{\raggedright\arraybackslash}p{0.12\columnwidth}@{}}
\toprule
Revisão & Âmbito & Estudos & Resposta automática & Onde e como atua & Avaliação da resposta \\
\midrule
Negera~\cite{negera2022review} (2022) & Deteção de \textit{botnets} com ML em IoT com SDN & n.i. & Não é o foco & Não & Não \\
Vera-Mundaca~\cite{veramundaca2024cybersecurity} (2024) & Ameaças e mitigação em casas inteligentes & 49 (2020--2024) & Parcial (estratégias) & Não & Não \\
Kandah~\cite{kandah2025navigating} (2025) & Arquiteturas, ameaças e contramedidas em IoT & n.i. & Parcial (contramedidas) & Não & Não \\
Bhardwaj~\cite{bhardwaj2025systematic} (2025) & \textit{Threat intelligence} e IDS (empresa, \textit{cloud}, IoT) & n.i. & Parcial (tendência) & Não & Não \\
\midrule
\textbf{Esta revisão} & Deteção e contenção de dispositivos IoT comprometidos com ML & 31 (2020--2026) & Sim (EC6, QF3) & Sim (QF3) & Sim (QF4, $T_c$) \\
\bottomrule
\end{tabular}
\end{table}



A leitura das Tabelas~\ref{tab:comparacao} e~\ref{tab:comparacao2}, com os 31 estudos incluídos, mostra quatro padrões. O primeiro é que a resposta é aplicada
sobretudo na infraestrutura de rede programável (\textit{switches} P4 ou controladores
SDN)~\cite{ramirezmartinez2025p4,andrade2024ddos,elsayed2024mpguard,zolotukhin2021curious,tawfik2021internet,Iqbal2026,Chaudhary2025,toony2024multiblock,gowtham2026fedmp,mazhar2022ridps,garba2024sdn,karibandi2025cyberdome},
o que pressupõe equipamento que uma rede doméstica ou de um pequeno edifício raramente
tem. O segundo é que poucos estudos medem a contenção em si: entre os que o fazem, os
tempos vão de milissegundos~\cite{toony2024multiblock,shukla2026cyber} e décimas de
segundo~\cite{ramirezmartinez2025p4,kponyo2020lightweight} a alguns
segundos~\cite{andrade2024ddos,karam2025isaaf,srilekha2025hybrid}, e os valores mais baixos
são obtidos em emulação ou simulação. O terceiro é que, em 14 dos 31 estudos, a
resposta não chega a ser avaliada numa rede: é apenas descrita, avaliada só
\textit{offline} sobre o \textit{dataset} ou anunciada sem ser
implementada~\cite{dange2025dual,tawfik2021internet,saurabh2022mlaps,rawat2024intelligent,Rasheed2026,Arulkumar2026,Chaudhary2025,shukla2026cyber,hussain2026postquantum,verma2024petrak,garba2024sdn,tariq2025saegru,kumar2025advanced,karibandi2025cyberdome}.
O quarto é que só sete trabalhos foram testados em \textit{hardware}
real~\cite{ravi2020semisupervised,dange2025dual,kponyo2020lightweight,Iqbal2026,karam2025isaaf,srilekha2025hybrid,garba2024sdn},
e os que medem o tempo de contenção nesse \textit{hardware} ou tratam apenas
\textit{floods} e decidem com base na CPU e na memória~\cite{kponyo2020lightweight}, ou
levam alguns segundos a reagir~\cite{karam2025isaaf,srilekha2025hybrid}.
Nenhum dos estudos comparados junta, num \textit{gateway} comum, a extração de fluxos, a
comparação entre modelos supervisionados e não supervisionados e uma quarentena
automática com o tempo de contenção medido. É esse o espaço que a solução proposta
procura ocupar.

% =====================================================================
\section{Solução proposta}\label{sec:proposta}
% Ponte para a Parte 2. RESPONSÁVEL: Frente 1 (arquitetura) + Frentes 3, 4, 5
% =====================================================================
A solução que se propõe para a Parte~2 junta, num único equipamento, as duas funções
que a revisão mostra estarem muitas vezes separadas: a deteção de dispositivos
comprometidos a partir do tráfego e a sua contenção automática. Tudo corre no
\textit{gateway} da rede IoT, que é o único ponto por onde passa todo o tráfego dos
dispositivos e onde é possível bloqueá-lo sem instalar nada nos equipamentos. A escolha
segue a direção apontada por Noetzold \textit{et al.}~\cite{noetzold2026anomaly}, que
deixam a resposta automática e a execução na \textit{edge} como trabalho futuro, e
aproxima-se dos estudos que aplicam a contenção junto à origem do tráfego, num nó
\textit{fog}~\cite{ravi2020semisupervised} ou num comutador de
fronteira~\cite{ramirezmartinez2025p4}.

\subsection{Arquitetura}
O cenário de teste é montado com contentores Docker num único computador, o que permite
repetir as experiências sem \textit{hardware} dedicado e sem pôr em risco outras redes
(Fig.~\ref{fig:arquitetura}). Há duas redes virtuais. Na rede IoT estão um ou dois
contentores que fazem de dispositivos e um contentor atacante; na rede externa está um
servidor que representa a Internet. O \textit{gateway} liga as duas redes,
encaminha o tráfego entre elas e concentra quatro módulos:
\begin{enumerate}
  \item \textbf{Captura e extração de fluxos.} O Zeek~\cite{paxson1999bro} observa a
  interface da rede IoT e escreve um registo por ligação (\texttt{conn.log}), com
  campos como a duração, os bytes e pacotes em cada sentido, o protocolo, o serviço, a
  porta de destino e o estado da ligação.
  \item \textbf{Deteção.} Um serviço em Python lê os registos à medida que são escritos,
  normaliza as \textit{features} e aplica os modelos treinados com
  scikit-learn~\cite{pedregosa2011scikit}.
  \item \textbf{Decisão.} Agrega as previsões por dispositivo, isto é, por endereço IP
  de origem, e decide se o dispositivo deve ser isolado (Secção~\ref{sec:decisao}).
  \item \textbf{Contenção.} Acrescenta o endereço do dispositivo a um conjunto de
  quarentena no nftables~\cite{nftables}, com um tempo de expiração. Uma regra na
  cadeia de encaminhamento descarta todo o tráfego cuja origem ou destino pertença a
  esse conjunto, e cada decisão fica registada com a hora e o motivo.
\end{enumerate}

Por razões éticas e de segurança, não é executado \textit{malware} real. O tráfego
malicioso vem das capturas do IoT-23~\cite{iot23}, que é reproduzido na rede IoT pelo
contentor atacante com o \texttt{tcpreplay}, depois de os endereços serem reescritos
com o \texttt{tcprewrite} para corresponderem aos dos contentores. Para o
\textit{gateway}, este tráfego é indistinguível do de um dispositivo infetado.

\begin{figure}[H]
\centering
\begin{tikzpicture}[font=\scriptsize, >={Stealth[length=1.5mm]},
  caixa/.style={draw, rounded corners=1pt, align=center, inner sep=2.5pt, minimum height=5mm},
  modulo/.style={draw, fill=gray!10, align=center, inner sep=2pt, minimum width=24mm, minimum height=4.5mm}]
  % rede IoT
  \node[caixa] (d1) at (0,1.15) {Dispositivo\\IoT 1};
  \node[caixa] (d2) at (0,0)    {Dispositivo\\IoT 2};
  \node[caixa] (at) at (0,-1.25) {Atacante\\(\texttt{tcpreplay}\\IoT-23)};
  \node[draw, dashed, fit=(d1)(d2)(at), inner sep=3pt, label={[font=\scriptsize]above:Rede IoT}] (riot) {};
  % gateway
  \node[modulo] (zeek) at (3.0,1.05)  {Zeek (\texttt{conn.log})};
  \node[modulo] (det)  at (3.0,0.35)  {Deteção (RF, IF, AE)};
  \node[modulo] (dec)  at (3.0,-0.35) {Decisão por dispositivo};
  \node[modulo] (nft)  at (3.0,-1.05) {nftables (quarentena)};
  \node[draw, thick, fit=(zeek)(det)(dec)(nft), inner sep=3pt, label={[font=\scriptsize]above:\textit{Gateway}}] (gw) {};
  \draw[->] (zeek) -- (det); \draw[->] (det) -- (dec); \draw[->] (dec) -- (nft);
  % rede externa
  \node[caixa] (srv) at (5.85,0) {Servidor\\externo};
  \node[draw, dashed, fit=(srv), inner sep=3pt, label={[font=\scriptsize]above:Rede externa}] {};
  % ligações
  \draw[<->] (riot.east) -- (gw.west);
  \draw[<->] (gw.east) -- (srv.west);
\end{tikzpicture}
\caption{Arquitetura do sistema proposto, em contentores Docker.}\label{fig:arquitetura}
\end{figure}

A escolha do IoT-23 tem uma vantagem prática importante: as capturas são distribuídas
com os registos \texttt{conn.log} do Zeek já rotulados. Os modelos são, por isso,
treinados com as mesmas \textit{features} que o \textit{gateway} calcula em
funcionamento. Evita-se assim um problema visível em alguns dos estudos analisados, que
treinam com as \textit{features} de um \textit{dataset} público e depois as tentam
reproduzir de outra forma no protótipo~\cite{elsayed2024mpguard}.

\subsection{Modelos de deteção}
% RESPONSÁVEL: Frentes 3 e 4
São comparados três modelos, um de cada família identificada na QF1. A ideia é avaliar o
compromisso entre detetar bem os ataques conhecidos e conseguir sinalizar
comportamentos que não apareceram no treino. As \textit{features} numéricas de cada
fluxo são normalizadas para $[0,1]$:
\begin{equation}
x' = \frac{x - x_{\min}}{x_{\max} - x_{\min}}
\label{eq:minmax}
\end{equation}
e as categóricas (protocolo, serviço e estado da ligação) são codificadas com
\textit{one-hot encoding}.

A \textit{Random Forest}~\cite{breiman2001random} é o modelo supervisionado. É treinada
com fluxos rotulados como benignos ou maliciosos e combina $B$ árvores de decisão,
cada uma treinada numa amostra \textit{bootstrap} dos dados e com um subconjunto
aleatório de \textit{features} em cada divisão. A probabilidade de um fluxo ser
malicioso é a fração de árvores que o classificam como tal:
\begin{equation}
p(x) = \frac{1}{B}\sum_{b=1}^{B} \mathbf{1}\left[h_b(x) = 1\right]
\label{eq:rf}
\end{equation}
É um dos modelos mais frequentes na deteção de intrusões em IoT e serve de referência
para os outros dois.

O \textit{Isolation Forest} é o modelo não supervisionado. É treinado sem rótulos e
classifica cada fluxo pelo \textit{score} de anomalia da Eq.~\ref{eq:iforest}.

O \textit{autoencoder}, implementado como uma rede MLP com uma camada central mais
estreita do que a entrada (Eqs.~\ref{eq:ae_enc} e~\ref{eq:ae_dec}), é o modelo
semi-supervisionado. É treinado só com fluxos benignos, e um fluxo é considerado anómalo
quando o erro de reconstrução da Eq.~\ref{eq:ae} ultrapassa o limiar $\theta$. O limiar $\theta$ é fixado num percentil alto do erro observado no tráfego benigno de
validação (por exemplo, o percentil 99), o que controla diretamente a taxa de falsos
positivos esperada.

Para evitar resultados otimistas, a divisão entre treino e teste é feita por captura e
não por fluxo: os fluxos de uma mesma captura do IoT-23 ficam todos do mesmo lado. Assim
o modelo nunca é testado com tráfego quase idêntico ao que viu no treino.

\subsection{Decisão e contenção}\label{sec:decisao}
Isolar um dispositivo por causa de um único fluxo suspeito seria arriscado, porque
mesmo uma taxa de falsos positivos baixa, multiplicada pelos milhares de fluxos que um
\textit{gateway} vê por hora, levaria a quarentenas indevidas frequentes. Os estudos
revistos que implementam a contenção resolvem isto exigindo várias evidências antes de
agir: um bloqueio só depois de três fluxos maliciosos~\cite{ramirezmartinez2025p4}, ou
de cinco previsões de ataque seguidas~\cite{andrade2024ddos}. Segue-se a mesma
abordagem. Para cada dispositivo $i$, considera-se a janela $W_i(t)$ com os seus
últimos $w$ fluxos, e o dispositivo é colocado em quarentena quando
\begin{equation}
D_i(t) = \mathbf{1}\left[\,\sum_{j \in W_i(t)} y_j \;\ge\; k\,\right],
\qquad y_j \in \{0,1\},
\label{eq:decisao}
\end{equation}
em que $y_j = 1$ se o fluxo $j$ foi sinalizado pelo modelo em uso. Os valores de $w$ e de
$k$ são escolhidos na Parte~2, procurando o melhor equilíbrio entre o tempo até à
contenção e o número de quarentenas indevidas.

A quarentena tem um tempo de expiração, tal como o bloqueio temporário usado por Ravi
e Shalinie~\cite{ravi2020semisupervised} e a quarentena de duração fixa de Ibrahim
\textit{et al.}~\cite{ibrahim2026purechain}. Se o dispositivo voltar a ser sinalizado
depois de libertado, a quarentena é renovada. Isto reduz o custo de um falso positivo
sem deixar um dispositivo realmente comprometido circular livremente. A regra bloqueia
o dispositivo inteiro e não apenas a ligação suspeita, à semelhança da regra global de
isolamento de \textit{host} do IDERES~\cite{rose2022ideres}, porque um equipamento
comprometido pode mudar de destino ou de porta a qualquer momento.

\subsection{Avaliação planeada}
% RESPONSÁVEL: Frente 1 (Afonso)
A avaliação responde às mesmas perguntas que foram usadas para analisar a literatura e
decorre em três cenários:
\begin{itemize}
  \item \textbf{S1, só tráfego benigno:} mede as quarentenas indevidas;
  \item \textbf{S2, ataque conhecido:} reprodução de capturas de famílias presentes no
  treino, como a Mirai;
  \item \textbf{S3, ataque não visto:} reprodução de capturas de uma família excluída do
  treino, para comparar a \textit{Random Forest} com os modelos que só aprendem o
  comportamento normal.
\end{itemize}

A deteção é avaliada \textit{offline} com as métricas que derivam da matriz de
confusão~\cite{sokolova2009systematic}, onde TP, FP, FN e TN são os verdadeiros positivos,
falsos positivos, falsos negativos e verdadeiros negativos:
\begin{equation}
\mathrm{Precision}=\frac{TP}{TP+FP},\qquad \mathrm{Recall}=\frac{TP}{TP+FN}
\label{eq:precrec}
\end{equation}
\begin{equation}
F_1 = \frac{2\cdot \mathrm{Precision}\cdot \mathrm{Recall}}{\mathrm{Precision}+\mathrm{Recall}}
\label{eq:f1}
\end{equation}
Num sistema que isola dispositivos, a taxa de falsos positivos é a que mede o custo de
isolar um dispositivo legítimo, e por isso é reportada à parte:
\begin{equation}
FPR = \frac{FP}{FP+TN}
\label{eq:fpr}
\end{equation}
A estas junta-se a \textit{Polygon Area Metric} (PAM)~\cite{aydemir2020pam}, também usada
em~\cite{noetzold2026anomaly}. A PAM representa $n$ métricas normalizadas num gráfico
radar e calcula a área do polígono que formam, dividida pela área máxima possível. A
proposta original usa seis métricas (\textit{accuracy}, sensibilidade, especificidade,
AUC, índice de Jaccard e F1):
\begin{equation}
PAM = \frac{\frac{1}{2}\sin\left(\frac{2\pi}{n}\right)\sum_{i=1}^{n} r_i\, r_{i+1}}
{\frac{n}{2}\sin\left(\frac{2\pi}{n}\right)}
= \frac{1}{n}\sum_{i=1}^{n} r_i\, r_{i+1},
\label{eq:pam}
\end{equation}
com $r_{n+1} = r_1$. Um valor de 1 corresponde a um classificador perfeito em todas as
métricas.

No protótipo, mede-se o tempo até à contenção, desde o primeiro pacote do ataque até à
entrada do dispositivo na quarentena:
\begin{equation}
T_c = t_{\text{isolamento}} - t_{\text{ataque}} = T_{\text{deteção}} + T_{\text{resposta}}
\label{eq:tc}
\end{equation}
Registam-se também a utilização de CPU e de memória do \textit{gateway}, para perceber se
a solução é viável num equipamento de \textit{edge} com poucos recursos. Estes dados
respondem, para a solução proposta, à mesma QF4 que foi colocada aos estudos revistos.

% =====================================================================
\section{Enquadramento na arquitetura Zero Trust}\label{sec:zta}
% RESPONSÁVEL: Frente 1 (Afonso) + par do laboratório ZTA
% =====================================================================
A solução da Secção~\ref{sec:proposta} nasceu do problema dos dispositivos IoT
comprometidos, mas as suas peças correspondem quase uma a uma aos componentes de uma
arquitetura \textit{Zero Trust} (ZTA). Esta secção faz esse enquadramento. Usa como
referência o NIST SP 800-207~\cite{rose2020zta}, o guia de implementação NIST SP
1800-35~\cite{nccoe2025zta} e o modelo de maturidade da CISA~\cite{cisa2023ztmm}, e
aplica-os a um cenário concreto: um edifício inteligente de um campus universitário,
inserido numa cidade digital.

\subsection{Porquê \textit{Zero Trust} neste cenário}
O modelo de segurança tradicional confia em tudo o que está dentro do perímetro. O
problema, como nota o SP 800-207, é que depois de um atacante passar o perímetro o
movimento lateral fica praticamente livre. A expressão \textit{zero trust} foi
popularizada por Kindervag~\cite{kindervag2010zero} precisamente para inverter esse
pressuposto: nenhum sujeito ou ativo é confiável só por estar na rede interna, e cada
acesso é decidido por pedido, com o menor privilégio possível, como se a rede já
estivesse comprometida~\cite{rose2020zta}. A CISA leva a ideia mais longe e parte do
princípio de que uma violação já ocorreu ou vai ocorrer~\cite{cisa2023ztmm}.

O edifício do cenário tem câmaras IP, sensores ambientais, fechaduras e controladores
de climatização e iluminação, muitos instalados e mantidos por empresas externas. É o
caso que o SP 800-207 descreve para prestadores e visitantes (§4.3), com sistemas de
climatização e iluminação geridos por terceiros. Hoje estes equipamentos partilham uma
rede plana com o resto do edifício. Um único dispositivo infetado por uma variante da
Mirai pode varrer a rede, tentar credenciais por omissão noutros equipamentos e
participar em ataques DDoS, sem que nada o trave. O próprio SP 800-207 refere a Mirai
quando discute a negação de serviço contra a infraestrutura de decisão (§5.2).

Há uma razão adicional para tratar estes dispositivos com \textit{Zero Trust}, e não só
com mais regras de \textit{firewall}: os dispositivos IoT quase não têm identidade
forte. He \textit{et al.}~\cite{he2022survey} sublinham que estes equipamentos têm poucos
atributos únicos e persistentes que permitam autenticá-los, e que os protocolos baseados
em certificados são pesados para eles. O SP 800-207 lembra ainda que o endereço MAC não
serve como identidade, porque pode ser falsificado (§3.4.1). Como não se pode confiar na
identidade declarada, a confiança tem de vir do que o dispositivo faz na rede.

Os fatores que levam o campus a adotar este modelo são de três tipos. Há a ameaça em si:
\textit{botnets} como a Mirai procuram ativamente equipamentos deste género. Há a
operação: os prestadores externos precisam de aceder remotamente aos seus sistemas para
manutenção, o que abre caminhos para dentro da rede. E há as obrigações legais e de
serviço: as câmaras e as fechaduras tratam dados pessoais, abrangidos pelo RGPD, e uma
falha nos controladores afeta o funcionamento do edifício. Os desafios também são
conhecidos. O SP 1800-35 aponta a falta de inventário de ativos, a falta de visibilidade
sobre os padrões de comunicação, a dificuldade de integrar tecnologias de maturidades
diferentes e o receio de prejudicar as operações~\cite{nccoe2025zta}. Neste cenário
juntam-se mais dois: não é possível instalar agentes nos dispositivos, e isolar por
engano uma fechadura ou uma câmara tem consequências físicas.

\subsection{Princípios e componentes}
O SP 800-207 define sete princípios (\textit{tenets}, §2.1). A
Tabela~\ref{tab:zta_principios} mostra como cada um é tratado na solução proposta.
Ela cobre bem os princípios ligados à monitorização contínua e à política dinâmica
(4, 5 e 7). Os princípios ligados à autenticação e à cifra (2, 3 e 6) ficam fora do
âmbito da Parte~2: dependem de capacidades que os dispositivos do cenário não têm.

\begin{table}[H]
\caption{Princípios do NIST SP 800-207 na solução proposta (\checkok{} coberto, $\sim$ parcial, \nao{} fora do âmbito)}\label{tab:zta_principios}
\centering\footnotesize
\begin{tabular}{@{}p{0.40\columnwidth}p{0.08\columnwidth}p{0.42\columnwidth}@{}}
\toprule
Princípio (§2.1) & & Na solução proposta \\
\midrule
1. Todas as fontes de dados e serviços são recursos & \checkok & Cada dispositivo IoT é tratado como recurso com política própria \\
2. Comunicação protegida, qualquer que seja a localização & \nao & Fora do âmbito: muitos IoT não suportam cifra robusta \\
3. Acesso concedido por sessão & $\sim$ & Cada fluxo é avaliado; o acesso pode ser cortado a meio \\
4. Acesso decidido por política dinâmica & \checkok & A política muda com o comportamento observado (Eq.~\ref{eq:decisao}) \\
5. Monitorizar a integridade e a postura de todos os ativos & \checkok & Detetor ML sobre os fluxos de cada dispositivo \\
6. Autenticação e autorização dinâmicas antes do acesso & $\sim$ & Autorização revista em contínuo; sem autenticação forte \\
7. Recolher informação e usá-la para melhorar a política & \checkok & Registo de fluxos e decisões; re-treino dos modelos \\
\bottomrule
\end{tabular}
\end{table}

Os componentes lógicos de uma ZTA (§3.1) são três. O \textit{Policy Engine} (PE)
decide se um acesso é concedido, negado ou revogado. O \textit{Policy Administrator}
(PA) executa essa decisão, configurando o caminho entre o sujeito e o recurso. O
\textit{Policy Enforcement Point} (PEP) abre, monitoriza e fecha as ligações. PE e PA
formam o \textit{Policy Decision Point} (PDP). O PE apoia-se em várias fontes de
informação, a que o SP 1800-35 chama \textit{Policy Information Points}
(PIP)~\cite{nccoe2025zta}: o estado dos ativos, os registos de atividade da rede, a
\textit{threat intelligence} e a análise de segurança, entre outras.

A decisão do PE é tomada pelo algoritmo de confiança (§3.3). O SP 800-207 distingue
algoritmos por critérios, em que todos os requisitos têm de ser cumpridos, de
algoritmos por pontuação, que calculam um nível de confiança a partir das várias fontes
e o comparam com um limiar. Distingue também algoritmos singulares, que avaliam cada
pedido isoladamente, de algoritmos contextuais, que usam o histórico recente do sujeito
e por isso detetam comportamentos atípicos. Em notação, um algoritmo por pontuação pode
escrever-se como
\begin{equation}
C = \sum_{m} w_m\, s_m, \qquad
\text{acesso} \iff C \le L_{\text{recurso}},
\label{eq:zta_score}
\end{equation}
em que $s_m$ é o valor dado pela fonte $m$, $w_m$ o peso que a organização lhe atribui e
$L_{\text{recurso}}$ o limiar aceite para o recurso. Aqui $C$ é escrito como nível de
suspeita, para coincidir com o sentido das saídas dos modelos.

\subsection{A solução proposta como ZTA}
A Tabela~\ref{tab:zta_componentes} e a Fig.~\ref{fig:zta} mostram a correspondência
entre os módulos do \textit{gateway} e os componentes do SP 800-207. Duas escolhas
merecem justificação.

A primeira é a variante de implementação. O SP 800-207 descreve quatro (§3.2):
agente e \textit{gateway}, enclave, portal de recursos e \textit{sandboxing}. A
variante com agente pressupõe software instalado no dispositivo, o que é impossível
numa câmara IP ou num controlador de climatização. Por isso a solução segue o modelo de
enclave: um \textit{gateway} na fronteira de um grupo de recursos que não consegue
proteger-se sozinho. A abordagem é a microssegmentação (§3.1.2). Cada classe de
dispositivos fica num segmento próprio, e todo o tráfego entre segmentos e para a
Internet atravessa o PEP.

A segunda é o algoritmo de confiança. O detetor é um algoritmo contextual por
pontuação. O contexto é o histórico recente de fluxos de cada dispositivo, e a
pontuação é a saída dos modelos: a probabilidade da \textit{Random Forest}
(Eq.~\ref{eq:rf}), o \textit{score} do \textit{Isolation Forest} (Eq.~\ref{eq:iforest})
ou o erro do \textit{Autoencoder} (Eq.~\ref{eq:ae}). Na versão da Parte~2 há uma só fonte
de informação. A Eq.~\ref{eq:zta_score} reduz-se então a $C_i(t)=\frac{1}{w}\sum_{j\in
W_i(t)} y_j$, e a quarentena é imposta quando $C_i(t) \ge k/w$, o que é exatamente a
regra da Eq.~\ref{eq:decisao}. A formulação geral deixa espaço para acrescentar outras
fontes com pesos próprios, como listas de servidores de C\&C conhecidos ou o inventário
de dispositivos.

\begin{table}[H]
\caption{Componentes do SP 800-207 e módulos da solução}\label{tab:zta_componentes}
\centering\footnotesize
\begin{tabular}{@{}p{0.25\columnwidth}p{0.66\columnwidth}@{}}
\toprule
Componente & Na solução proposta \\
\midrule
Sujeitos e recursos & Dispositivos IoT do edifício (identidades de máquina) e serviços que contactam \\
PIP / fontes de dados & Zeek (registos de fluxos), detetor ML, inventário de dispositivos; \textit{threat intelligence} como extensão \\
PE & Módulo de decisão por dispositivo: algoritmo contextual por pontuação (Eq.~\ref{eq:decisao}) \\
PA & Serviço que traduz a decisão numa operação sobre o conjunto de quarentena \\
PEP & \textit{Gateway} com nftables, modelo de enclave por segmento IoT \\
Plano de controlo & Comunicação local entre detetor, decisão e nftables, separada do tráfego dos dispositivos \\
\bottomrule
\end{tabular}
\end{table}

\begin{figure}[H]
\centering
\begin{tikzpicture}[font=\scriptsize, >={Stealth[length=1.5mm]},
  bloco/.style={draw, rounded corners=1pt, align=center, inner sep=2.5pt, minimum height=5mm},
  pdp/.style={draw, fill=gray!12, align=center, inner sep=2.5pt, minimum height=5mm}]
  % plano de controlo
  \node[pdp] (pe) at (0,2.2) {PE\\decisão por dispositivo};
  \node[pdp] (pa) at (2.6,2.2) {PA\\gestão da quarentena};
  \node[draw, dashed, fit=(pe)(pa), inner sep=3pt, label={[font=\scriptsize]above:PDP (plano de controlo)}] (pdpbox) {};
  \node[bloco] (pip) at (-2.35,2.2) {PIP\\Zeek + ML\\inventário};
  % plano de dados
  \node[bloco] (iot) at (-2.35,0) {Segmento IoT\\(câmaras, AVAC,\\sensores)};
  \node[draw, thick, fill=gray!25, align=center, inner sep=3pt, minimum height=7mm] (pep) at (0.9,0) {PEP\\\textit{gateway} + nftables};
  \node[bloco] (rec) at (3.35,0) {Recursos e\\Internet};
  \draw[<->, thick] (iot) -- node[below, font=\tiny]{plano de dados} (pep);
  \draw[<->, thick] (pep) -- (rec);
  % ligações de controlo
  \draw[->] (pip) -- (pe);
  \draw[->] (pe) -- (pa);
  \draw[->, dashed] (pa.south) -- node[right, font=\tiny]{regra} (pep.north east);
  \draw[->, dotted] (pep.north west) -- node[left, font=\tiny, pos=0.35]{metadados} (pip.south east);
\end{tikzpicture}
\caption{A solução proposta segundo os componentes lógicos do NIST SP 800-207.}\label{fig:zta}
\end{figure}

A Tabela~\ref{tab:zta_casos} descreve os casos de uso, alinhados com os do SP
1800-35~\cite{nccoe2025zta}. O guia do NCCoE demonstra, entre outros, a descoberta de
ativos (caso~A) e a revisão contínua da confiança durante a sessão (casos~F). Nestes
últimos, um \textit{endpoint} que deixa de estar conforme perde o acesso, recupera-o
quando volta à conformidade, e um \textit{endpoint} suspeito de estar comprometido vê os
seus pedidos negados e as sessões terminadas. A diferença é que, no guia, a conformidade
é verificada por agentes instalados nos equipamentos; aqui é inferida do tráfego.

\begin{table}[H]
\caption{Casos de uso da solução}\label{tab:zta_casos}
\centering\footnotesize
\begin{tabular}{@{}p{0.04\columnwidth}p{0.62\columnwidth}p{0.24\columnwidth}@{}}
\toprule
& Situação e resposta & SP 1800-35 \\
\midrule
1 & Um dispositivo novo liga-se: é registado no inventário e o seu tráfego passa a ser observado & A (descoberta) \\
2 & Um dispositivo começa a varrer a rede ou a contactar um C\&C: é colocado em quarentena e todas as ligações são cortadas & F-14 a F-17 \\
3 & A quarentena expira e o tráfego volta ao normal: o dispositivo recupera o acesso; se reincidir, a quarentena é renovada & F-4 e F-5 \\
4 & Um dispositivo crítico é isolado por engano: o administrador liberta-o manualmente e o caso entra no re-treino & -- \\
\bottomrule
\end{tabular}
\end{table}

\subsection{Maturidade e plano de transição}
O modelo da CISA~\cite{cisa2023ztmm} organiza a ZTA em cinco pilares (identidade,
dispositivos, redes, aplicações e dados), três capacidades transversais (visibilidade e
análise, automação e orquestração, governação) e quatro níveis de maturidade
(tradicional, inicial, avançado e ótimo). Nos pilares que importam aqui, dispositivos e
redes, o edifício parte do nível tradicional: rede plana, inventário manual e nenhuma
resposta automática. A solução proposta leva a visibilidade ao nível avançado, que a
CISA associa à deteção de anomalias automática na rede. Aproxima a automação do nível
ótimo, que prevê processos automáticos para monitorizar, isolar e remediar dispositivos.
Ali \textit{et al.}~\cite{ali2022maturity} resumem bem este passo ao afirmar que a função
central da automação é converter a deteção de ameaças em prevenção.

A Tabela~\ref{tab:zta_cisa} resume as lacunas da situação de partida e o nível que a
solução permite atingir em cada função relevante. O pilar da identidade fica limitado ao
acesso dos administradores ao PDP, que deve exigir autenticação multifator. Para os
dispositivos, como se viu, a identidade forte não é realista.

\begin{table}[H]
\caption{Lacunas da situação de partida e nível atingido (CISA ZTMM v2)}\label{tab:zta_cisa}
\centering\footnotesize
\begin{tabular}{@{}>{\raggedright\arraybackslash}p{0.22\columnwidth}>{\raggedright\arraybackslash}p{0.32\columnwidth}>{\raggedright\arraybackslash}p{0.36\columnwidth}@{}}
\toprule
Pilar / função & Situação de partida & Com a solução proposta \\
\midrule
Dispositivos: visibilidade e análise & Inventário manual, sem deteção (tradicional) & Inventário a partir dos fluxos e deteção automática de anomalias (avançado) \\
Dispositivos: automação & Isolamento manual, quando existe (tradicional) & Quarentena automática com expiração e renovação (entre avançado e ótimo) \\
Redes: segmentação & Rede plana (tradicional) & Microssegmentação por classe, com o \textit{gateway} como PEP (avançado) \\
Redes: visibilidade e análise & Monitorização só na fronteira (tradicional) & Deteção de anomalias com ML em todo o tráfego IoT (avançado) \\
Identidade & Contas locais sem MFA (tradicional) & MFA no acesso dos administradores ao PDP (inicial) \\
\bottomrule
\end{tabular}
\end{table}

A transição segue os passos de migração do SP 800-207 (§7.3) e a sequência proposta pelo
SP 1800-35:
\begin{enumerate}
  \item inventariar os dispositivos e os fluxos de que dependem, com o próprio Zeek;
  \item segmentar a rede por classe de dispositivo, com o \textit{gateway} como único
  caminho entre segmentos;
  \item pôr o detetor em modo só de relatório, como recomenda o §7.3.6, para conhecer o
  comportamento normal e afinar $w$ e $k$ sem afetar o serviço;
  \item ativar a quarentena automática, primeiro nos segmentos menos críticos;
  \item rever periodicamente os modelos e alargar as fontes de informação do PE.
\end{enumerate}

\subsection{Riscos da própria arquitetura}
Uma ZTA também introduz riscos, e o SP 800-207 dedica-lhes uma secção (§5). Três são
diretamente relevantes. O primeiro é a subversão do processo de decisão (§5.1): quem
controlar o detetor ou as regras pode isolar dispositivos legítimos ou libertar os
comprometidos. Mitiga-se com um plano de controlo isolado, acesso dos
administradores com autenticação multifator e registo de todas as decisões. O segundo é a negação de serviço contra
o PEP (§5.2): se o \textit{gateway} cair, o edifício perde a ligação, pelo que em
produção seria necessária redundância. O terceiro é o risco das entidades não pessoais
que tomam decisões, como modelos de ML (§5.7), cujo efeito mais visível são os falsos
positivos. A janela de $w$ fluxos, a quarentena com expiração e o período inicial só de
relatório respondem diretamente a este risco.

Há ainda uma questão de privacidade. Monitorizar o tráfego de um edifício pode tocar em
dados pessoais. Por isso a solução usa apenas metadados de fluxos, sem conteúdo,
seguindo o princípio da minimização do RGPD. O próprio SP 800-207 reconhece, no §5.4, que
técnicas de ML aplicadas a metadados permitem analisar tráfego que não pode ser
decifrado.

\subsection{Lacunas e contributo}
As referências oficiais deixam o caso IoT em aberto. O SP 1800-35 exclui explicitamente
do seu âmbito os sistemas de controlo industrial, a tecnologia operacional e os
dispositivos IoT~\cite{nccoe2025zta}. O modelo da CISA declara que não trata os desafios
de algumas classes de dispositivos IoT nem a incorporação de ML e IA em soluções
\textit{Zero Trust}~\cite{cisa2023ztmm}. Na literatura científica, Syed
\textit{et al.}~\cite{syed2022zta} apontam o cálculo da confiança e do risco como um dos
obstáculos a uma ZTA completa. He \textit{et al.}~\cite{he2022survey} apresentam o ML
como o método de avaliação de confiança mais preciso, mas também o mais exigente em
computação. Li \textit{et al.}~\cite{li2024iiot} propõem uma arquitetura
\textit{Zero Trust} para IoT industrial centrada na autenticação dos dispositivos, com
\textit{blockchain}.

Do lado da aplicação, o FURZE de Vanickis \textit{et al.}~\cite{vanickis2018access} é o
trabalho mais próximo: gera e retira regras de \textit{firewall} entre zonas à medida que
o risco de uma sessão muda. O risco, porém, é calculado com lógica difusa sobre
atributos, e não a partir do tráfego. O eZTrust de Zaheer \textit{et al.}~\cite{zaheer2019eztrust}
mostra que até entre microsserviços o endereço IP não é uma identidade fiável, mas exige
controlar o \textit{kernel} do emissor, o que não é possível num dispositivo IoT. A
abordagem MUD~\cite{rfc8520,nist2021mud} é complementar: o fabricante declara o
comportamento de rede esperado do dispositivo, e o \textit{gateway} aplica essa lista.
Depende, contudo, de os fabricantes publicarem estes perfis.

A solução proposta ocupa o espaço entre estes trabalhos. A confiança é inferida do
comportamento observado, sem exigir nada ao dispositivo nem ao fabricante. A inferência
corre no \textit{gateway}, o que contorna a falta de recursos dos dispositivos. E a
resposta é aplicada automaticamente pelo PEP. Isto dá aos dispositivos IoT a verificação
periódica da saúde dos recursos que o SP 1800-35 considera muito desejável e que vários
dos seus \textit{builds} não conseguiram assegurar~\cite{nccoe2025zta}.

% =====================================================================
\section{Discussão}\label{sec:discussao}
% RESPONSÁVEL: Frente 1 (Afonso); cada frente pode acrescentar à sua parte.
% =====================================================================
\subsection{Lacunas e oportunidades}
Da comparação resultam cinco lacunas, resumidas na Tabela~\ref{tab:lacunas} com a forma
como a solução proposta lhes responde.

A primeira é que a contenção é anunciada com mais frequência do que é avaliada. Mesmo
entre os estudos incluídos, que já passaram um critério de exclusão sobre a resposta
(EC6), vários descrevem ou anunciam um mecanismo que depois não avaliam numa
rede~\cite{elsayed2024mpguard,ibrahim2026purechain,dange2025dual,tawfik2021internet,saurabh2022mlaps,rawat2024intelligent},
e só sete medem quanto tempo demora a conter um
ataque~\cite{ramirezmartinez2025p4,andrade2024ddos,kponyo2020lightweight,karam2025isaaf,toony2024multiblock,srilekha2025hybrid,shukla2026cyber}.
A segunda é a dependência de infraestrutura programável. Boa parte das propostas
pressupõe \textit{switches} P4 ou um controlador SDN, que não existem numa casa, numa
pequena empresa ou num edifício com uma rede convencional, onde o ponto de controlo
natural é um \textit{gateway} Linux.

A terceira diz respeito aos dados. Os \textit{datasets} estão muito dispersos, vários
não são de IoT, e o tráfego próprio quase nunca é publicado, o que impede comparar
resultados entre estudos. A quarta é a escassa atenção aos falsos positivos: poucos
estudos reportam a FPR, embora seja a métrica que determina se a contenção automática é
aceitável. A resposta que a literatura encontra na prática, exigir evidência acumulada
antes de agir~\cite{ramirezmartinez2025p4,andrade2024ddos}, raramente é discutida
como tal. A quinta é que nenhum estudo compara de forma sistemática modelos
supervisionados e não supervisionados quanto ao efeito na contenção, embora a escolha
entre os dois determine se o sistema reage a ataques que não estavam no treino.

A estas juntam-se as lacunas das referências oficiais de \textit{Zero Trust}, que
deixam o IoT e o ML fora do seu âmbito~\cite{nccoe2025zta,cisa2023ztmm}, como se discutiu
na Secção~\ref{sec:zta}.

\begin{table}[H]
\caption{Lacunas identificadas e resposta da solução proposta}\label{tab:lacunas}
\centering\footnotesize
\begin{tabular}{@{}>{\raggedright\arraybackslash}p{0.43\columnwidth}>{\raggedright\arraybackslash}p{0.49\columnwidth}@{}}
\toprule
Lacuna & Resposta na solução proposta \\
\midrule
Contenção descrita mas pouco avaliada & Medição de $T_c$, separando deteção e resposta (Eq.~\ref{eq:tc}) \\
Dependência de P4 ou SDN & Contenção num \textit{gateway} comum, com nftables \\
\textit{Datasets} dispersos e não IoT & IoT-23, com os mesmos registos Zeek no treino e no protótipo \\
FPR pouco reportada & FPR, quarentenas indevidas (cenário S1) e decisão por $k$ de $w$ fluxos (Eq.~\ref{eq:decisao}) \\
Falta de comparação entre famílias de modelos & RF, IF e AE avaliados nos mesmos cenários, incluindo ataques não vistos (S3) \\
IoT e ML fora das referências de \textit{Zero Trust} & Enquadramento no NIST SP 800-207 e na CISA (Secção~\ref{sec:zta}) \\
\bottomrule
\end{tabular}
\end{table}

\subsection{Ameaças à validade}
Como em qualquer revisão sistemática, os resultados dependem das escolhas do
protocolo. Identificam-se quatro ameaças.

\textbf{Pesquisa.} A \textit{string} foi aplicada apenas ao resumo e exige termos ligados
à resposta. Estudos que só descrevem a contenção no texto completo podem ter ficado de
fora. No ScienceDirect foi usada uma versão reduzida da \textit{string}, e a pesquisa
limitou-se a cinco bases, a artigos em inglês e ao período 2020--2026.

\textbf{Seleção.} A triagem foi dividida por blocos, cada um com um revisor. Para reduzir
a subjetividade, os critérios foram escritos antes da triagem, os casos duvidosos
passaram sempre à etapa seguinte e as decisões ficaram registadas no Parsifal. O critério
EC6 foi introduzido depois da primeira leitura dos resumos. Embora seja um refinamento
previsto no protocolo e esteja documentado, pode ter afastado estudos com contributos
parciais sobre a resposta. Como cada bloco aplicou os critérios de forma independente, o
grau de exigência com que o EC6 foi lido também variou: alguns estudos foram incluídos
por proporem uma resposta, mesmo sem a avaliarem. Essa diferença fica visível na
Tabela~\ref{tab:comparacao} e foi tida em conta na leitura dos resultados. Na revisão
final do conjunto, seis estudos que tinham passado a avaliação de qualidade foram
excluídos pelo EC6, por não proporem qualquer mecanismo de resposta: limitavam-se à
classificação, a um alerta ou a uma notificação ao utilizador.
Cinco estudos foram triados por dois revisores, por aparecerem em dois blocos: em
quatro as decisões e as pontuações coincidiram, e no quinto~\cite{karibandi2025cyberdome}
um revisor excluiu-o (4,5) e outro incluiu-o (5,0), tendo prevalecido a inclusão.

\textbf{Extração e qualidade.} A avaliação de qualidade e a extração foram feitas por um
revisor por estudo. Em alguns artigos as tabelas e figuras não eram legíveis no texto
disponível, e vários contêm incoerências numéricas, como se referiu na
Secção~\ref{sec:resultados}. Os valores apresentados são os reportados pelos autores e não
foram reproduzidos.

\textbf{Generalização.} O número de estudos incluídos é reduzido, pelo que as proporções
indicadas descrevem este conjunto e não devem ser lidas como estimativas da área no seu
todo.

% =====================================================================
\section{Conclusões e trabalhos futuros}\label{sec:conclusoes}
% RESPONSÁVEL: Frente 1 (Afonso)
% =====================================================================
Esta revisão procurou saber como é que a análise de tráfego com ML é usada para detetar
e conter automaticamente dispositivos IoT comprometidos (QG). A resposta curta é que a
deteção está madura, mas a contenção ainda não. Dos 491 estudos únicos devolvidos pela
pesquisa, 31 cumpriram os critérios de seleção e a avaliação de qualidade. A maioria
deteta os ataques com uma exatidão acima de 97\%, mas só 14 descrevem e avaliam o
mecanismo de resposta, e só sete medem quanto tempo demora a conter um ataque.

Quanto aos modelos e às \textit{features} (QF1), domina a aprendizagem supervisionada: a
RF, as árvores de decisão, o KNN e a SVM aparecem ao lado de redes CNN, LSTM e GRU,
muitas vezes combinadas em \textit{ensembles}. As abordagens não supervisionadas e
semi-supervisionadas são poucas e, mesmo quando surgem ao lado de modelos
supervisionados, nenhum estudo mede o efeito de cada família na contenção. As
\textit{features} são, na grande maioria, estatísticas de fluxo, o que permite trabalhar
com tráfego cifrado e sem tocar nos dispositivos.

Quanto aos dados (QF2), os 31 estudos usam 29 \textit{datasets} públicos diferentes, 13
recorrem a pelo menos um que não é de IoT, e os dados próprios raramente são publicados.
O DoS/DDoS aparece em 28 estudos e é quase sempre o ataque contra o qual a resposta é
demonstrada. A FPR, que indica quantos dispositivos legítimos seriam isolados por engano,
só tem valor explícito em 13 estudos, e apenas nove medem a eficácia da própria resposta.

Quanto à resposta (QF3), o descarte ou o bloqueio de tráfego é a medida mais comum, e 11
estudos isolam o dispositivo inteiro. A resposta é aplicada sobretudo em controladores
SDN ou \textit{switches} P4 (12 estudos) e em \textit{gateways} ou nós \textit{edge} (12).
Os estudos mais cuidadosos não reagem a um único alerta: exigem evidência acumulada ou
usam modelos de decisão, como a aprendizagem por reforço. Em quatro, pelo contrário, a
resposta é proposta mas não chega a ser implementada.

Quanto à avaliação (QF4), sete estudos são avaliados apenas \textit{offline} e só sete
correm em \textit{hardware} real. Na maioria dos casos em que o tempo de contenção é
medido, a deteção demora uma a duas ordens de grandeza mais do que a resposta, dezenas de
segundos contra décimas de segundo. O custo computacional é reportado de forma desigual e
raramente no equipamento onde a deteção corre.

A comparação das Tabelas~\ref{tab:comparacao} e~\ref{tab:comparacao2} mostra a lacuna
que motiva este trabalho: nenhum estudo junta, num \textit{gateway} comum, a extração de
fluxos, a comparação entre modelos supervisionados e não supervisionados e uma quarentena
automática com o tempo de contenção medido. A solução proposta
(Secção~\ref{sec:proposta}) ocupa esse espaço com um \textit{gateway} em contentores que
extrai fluxos com o Zeek, compara a RF, o \textit{Isolation Forest} e o
\textit{autoencoder} nos mesmos cenários e isola com nftables os dispositivos com $k$
fluxos maliciosos numa janela de $w$. Enquadrada no NIST SP 800-207
(Secção~\ref{sec:zta}), a solução faz do \textit{gateway} um PEP e do detetor um
algoritmo de confiança contextual, e cobre um caso que as referências oficiais de
\textit{Zero Trust} deixam em aberto: o de dispositivos que não conseguem autenticar-se
nem correr agentes de segurança.

\textbf{Trabalhos futuros.} O passo imediato é a Parte~2: implementar o protótipo e
avaliá-lo nos cenários S1 a S3, com as métricas das
Eqs.~(\ref{eq:precrec})--(\ref{eq:fpr}), a PAM (Eq.~\ref{eq:pam}) e o tempo de contenção
$T_c$ (Eq.~\ref{eq:tc}), separado em deteção e resposta, a par da CPU e da memória do
\textit{gateway}. Os resultados dirão se os modelos que só aprendem o comportamento
normal compensam, perante ataques não vistos no treino, a menor precisão que costumam ter
face à RF. Depois disso, há quatro extensões naturais: uma resposta graduada, que limite
o débito antes de isolar e só isole quando a evidência persiste; mais fontes de
informação para o PE, como os perfis MUD dos fabricantes e o inventário de dispositivos;
a validação com outros \textit{datasets} e em \textit{hardware} real, com um Raspberry Pi
a fazer de \textit{gateway}; e a atualização contínua dos modelos, com aprendizagem
incremental, para acompanhar a mudança do tráfego legítimo ao longo do tempo.

% ---------------------------------------------------------------------
\bibliographystyle{IEEEtran}
\bibliography{referencias}

\end{document}
